\subsection{Additional Band-Composition and Region Galleries}
\label{sec:appendix_extra_region_gallery}

\captionsetup[figure]{font=small}
\captionsetup[subfigure]{font=small}

The figures below extend the staged-computation analysis from the main text through full region-wise curves rather than through band summaries. The retained \texttt{expr\_first} galleries show the entire one-hop/two-hop profile for each region and branch, which makes two points visible at layer resolution: prompt reordering moves Expression earlier, and Query remains the main integration sink even after that reordering.

\paragraph{How to read these region figures.}
Each panel should be read as a layer-wise sensitivity profile for one semantic region. A sharp early peak indicates a concentrated front-of-network computation stage; a broader middle or late plateau indicates that the same region remains active deeper into the network. The most informative comparison is within a branch: Facts versus Expression versus Query under the same model and prompt order.

%\begin{figure*}[p]
%    \centering
%    \begin{subfigure}[b]{0.29\textwidth}
%        \centering
%        \includegraphics[width=\linewidth]{./figures_experiments/reports/mlp_analysis/comparison/bcr_stacked.png}
%        \caption{\textit{Qwen3-8B}: MLP}
%    \end{subfigure}
%    \hfill
%    \begin{subfigure}[b]{0.29\textwidth}
%        \centering
%        \includegraphics[width=\linewidth]{./figures_experiments/reports/attn_analysis/comparison/bcr_stacked.png}
%        \caption{\textit{Qwen3-8B}: Attention}
%    \end{subfigure}
%    \hfill
%    \begin{subfigure}[b]{0.29\textwidth}
%        \centering
%        \includegraphics[width=\linewidth]{./figures_experiments/reports/attn_mlp_analysis/comparison/bcr_stacked.png}
%        \caption{\textit{Qwen3-8B}: Attention+MLP}
%    \end{subfigure}
%
%    \vspace{0.2em}
%
%    \begin{subfigure}[b]{0.29\textwidth}
%        \centering
%        \includegraphics[width=\linewidth]{./figures_experiments/reports/mlp_analysis/comparison_14B/bcr_stacked.png}
%        \caption{\textit{Qwen3-14B}: MLP}
%    \end{subfigure}
%    \hfill
%    \begin{subfigure}[b]{0.29\textwidth}
%        \centering
%        \includegraphics[width=\linewidth]{./figures_experiments/reports/attn_analysis/comparison_14B/bcr_stacked.png}
%        \caption{\textit{Qwen3-14B}: Attention}
%    \end{subfigure}
%    \hfill
%    \begin{subfigure}[b]{0.29\textwidth}
%        \centering
%        \includegraphics[width=\linewidth]{./figures_experiments/reports/attn_mlp_analysis/comparison_14B/bcr_stacked.png}
%        \caption{\textit{Qwen3-14B}: Attention+MLP}
%    \end{subfigure}
%    \caption{\textbf{Qwen3 \texttt{facts\_first}: stacked BCR summaries across scales and branches.} Each panel shows the early/middle/late contribution decomposition within the Facts, Expression, and Query regions. The MLP and joint branches maintain the clearest early-Facts profile, while the attention branch spreads more mass into middle-layer Expression and Query routing. The 14B results preserve the same ordering while reducing some of the 8B late-query mass.}
%    \label{fig:qwen3_bcr_stacked_facts}
%\end{figure*}

%\begin{figure*}[p]
%    \centering
%    \begin{subfigure}[b]{0.29\textwidth}
%        \centering
%        \includegraphics[width=\linewidth]{./figures_experiments/reports/mlp_analysis_expr_first/comparison_expr_first/bcr_stacked.png}
%        \caption{\textit{Qwen3-8B}: MLP}
%    \end{subfigure}
%    \hfill
%    \begin{subfigure}[b]{0.29\textwidth}
%        \centering
%        \includegraphics[width=\linewidth]{./figures_experiments/reports/attn_analysis_expr_first/comparison_expr_first/bcr_stacked.png}
%        \caption{\textit{Qwen3-8B}: Attention}
%    \end{subfigure}
%    \hfill
%    \begin{subfigure}[b]{0.29\textwidth}
%        \centering
%        \includegraphics[width=\linewidth]{./figures_experiments/reports/attn_mlp_analysis_expr_first/comparison_expr_first/bcr_stacked.png}
%        \caption{\textit{Qwen3-8B}: Attention+MLP}
%    \end{subfigure}
%
%    \vspace{0.2em}
%
%    \begin{subfigure}[b]{0.29\textwidth}
%        \centering
%        \includegraphics[width=\linewidth]{./figures_experiments/reports/mlp_analysis_expr_first/comparison_expr_first_14B/bcr_stacked.png}
%        \caption{\textit{Qwen3-14B}: MLP}
%    \end{subfigure}
%    \hfill
%    \begin{subfigure}[b]{0.29\textwidth}
%        \centering
%        \includegraphics[width=\linewidth]{./figures_experiments/reports/attn_analysis_expr_first/comparison_expr_first_14B/bcr_stacked.png}
%        \caption{\textit{Qwen3-14B}: Attention}
%    \end{subfigure}
%    \hfill
%    \begin{subfigure}[b]{0.29\textwidth}
%        \centering
%        \includegraphics[width=\linewidth]{./figures_experiments/reports/attn_mlp_analysis_expr_first/comparison_expr_first_14B/bcr_stacked.png}
%        \caption{\textit{Qwen3-14B}: Attention+MLP}
%    \end{subfigure}
%    \caption{\textbf{Qwen3 \texttt{expr\_first}: stacked BCR summaries across scales and branches.} Compared with Figure~\ref{fig:qwen3_bcr_stacked_facts}, the Expression Region acquires much more early-band mass, while the Facts Region becomes less exclusively early-layer concentrated. The effect is strongest in the attention and joint branches, but every branch still preserves a distinct query-facing aggregation component rather than collapsing into a single early computation stage.}
%    \label{fig:qwen3_bcr_stacked_expr}
%\end{figure*}

\begin{figure*}[t]
    \centering
    \begin{subfigure}[b]{0.285\textwidth}
        \centering
        \includegraphics[width=\linewidth]{./figures_experiments/reports/mlp_analysis_expr_first/mlp_8b_one_hop_expr_first/facts_region/mlp_facts_region_scores.png}
        \caption{Facts (One-hop)}
    \end{subfigure}
    \hfill
    \begin{subfigure}[b]{0.285\textwidth}
        \centering
        \includegraphics[width=\linewidth]{./figures_experiments/reports/mlp_analysis_expr_first/mlp_8b_one_hop_expr_first/expression_region/mlp_expression_region_scores.png}
        \caption{Expression (One-hop)}
    \end{subfigure}
    \hfill
    \begin{subfigure}[b]{0.285\textwidth}
        \centering
        \includegraphics[width=\linewidth]{./figures_experiments/reports/mlp_analysis_expr_first/mlp_8b_one_hop_expr_first/query_region/mlp_query_region_scores.png}
        \caption{Query (One-hop)}
    \end{subfigure}

    \vspace{0.2em}

    \begin{subfigure}[b]{0.285\textwidth}
        \centering
        \includegraphics[width=\linewidth]{./figures_experiments/reports/mlp_analysis_expr_first/mlp_8b_two_hop_expr_first/facts_region/mlp_facts_region_scores.png}
        \caption{Facts (Two-hop)}
    \end{subfigure}
    \hfill
    \begin{subfigure}[b]{0.285\textwidth}
        \centering
        \includegraphics[width=\linewidth]{./figures_experiments/reports/mlp_analysis_expr_first/mlp_8b_two_hop_expr_first/expression_region/mlp_expression_region_scores.png}
        \caption{Expression (Two-hop)}
    \end{subfigure}
    \hfill
    \begin{subfigure}[b]{0.285\textwidth}
        \centering
        \includegraphics[width=\linewidth]{./figures_experiments/reports/mlp_analysis_expr_first/mlp_8b_two_hop_expr_first/query_region/mlp_query_region_scores.png}
        \caption{Query (Two-hop)}
    \end{subfigure}
    \caption{\textbf{Complete \textit{Qwen3-8B} MLP region gallery under \texttt{expr\_first}.} The full one-hop/two-hop curves show that reordering the prompt moves expression processing earlier without removing the later query-facing integration stage. In one-hop, the Facts Region no longer dominates the earliest layers as strongly as under \texttt{facts\_first}, while the Expression Region becomes the sharpest early-layer signal. In two-hop, the same ordering survives but the Query Region broadens into a stronger middle-to-late bridge.}
    \label{fig:qwen3_8b_mlp_expr_gallery}
\end{figure*}

\begin{figure*}[t]
    \centering
    \begin{subfigure}[b]{0.285\textwidth}
        \centering
        \includegraphics[width=\linewidth]{./figures_experiments/reports/mlp_analysis_expr_first/mlp_14b_one_hop_expr_first/facts_region/mlp_facts_region_scores.png}
        \caption{Facts (One-hop)}
    \end{subfigure}
    \hfill
    \begin{subfigure}[b]{0.285\textwidth}
        \centering
        \includegraphics[width=\linewidth]{./figures_experiments/reports/mlp_analysis_expr_first/mlp_14b_one_hop_expr_first/expression_region/mlp_expression_region_scores.png}
        \caption{Expression (One-hop)}
    \end{subfigure}
    \hfill
    \begin{subfigure}[b]{0.285\textwidth}
        \centering
        \includegraphics[width=\linewidth]{./figures_experiments/reports/mlp_analysis_expr_first/mlp_14b_one_hop_expr_first/query_region/mlp_query_region_scores.png}
        \caption{Query (One-hop)}
    \end{subfigure}

    \vspace{0.2em}

    \begin{subfigure}[b]{0.285\textwidth}
        \centering
        \includegraphics[width=\linewidth]{./figures_experiments/reports/mlp_analysis_expr_first/mlp_14b_two_hop_expr_first/facts_region/mlp_facts_region_scores.png}
        \caption{Facts (Two-hop)}
    \end{subfigure}
    \hfill
    \begin{subfigure}[b]{0.285\textwidth}
        \centering
        \includegraphics[width=\linewidth]{./figures_experiments/reports/mlp_analysis_expr_first/mlp_14b_two_hop_expr_first/expression_region/mlp_expression_region_scores.png}
        \caption{Expression (Two-hop)}
    \end{subfigure}
    \hfill
    \begin{subfigure}[b]{0.285\textwidth}
        \centering
        \includegraphics[width=\linewidth]{./figures_experiments/reports/mlp_analysis_expr_first/mlp_14b_two_hop_expr_first/query_region/mlp_query_region_scores.png}
        \caption{Query (Two-hop)}
    \end{subfigure}
    \caption{\textbf{Complete \textit{Qwen3-14B} MLP region gallery under \texttt{expr\_first}.} The larger model exhibits the same prompt-order effect as 8B, but more cleanly: one-hop Expression is the most concentrated early-layer signal, Facts become broader and less purely early-band dominated, and Query remains the terminal aggregation region. Two-hop again lowers the absolute sharpness of every region while preserving the same coarse temporal ordering.}
    \label{fig:qwen3_14b_mlp_expr_gallery}
\end{figure*}

\begin{figure*}[t]
    \centering
    \begin{subfigure}[b]{0.285\textwidth}
        \centering
        \includegraphics[width=\linewidth]{./figures_experiments/reports/attn_analysis_expr_first/attn_8b_one_hop_expr_first/facts_region/attn_facts_region_scores.png}
        \caption{Facts (One-hop)}
    \end{subfigure}
    \hfill
    \begin{subfigure}[b]{0.285\textwidth}
        \centering
        \includegraphics[width=\linewidth]{./figures_experiments/reports/attn_analysis_expr_first/attn_8b_one_hop_expr_first/expression_region/attn_expression_region_scores.png}
        \caption{Expression (One-hop)}
    \end{subfigure}
    \hfill
    \begin{subfigure}[b]{0.285\textwidth}
        \centering
        \includegraphics[width=\linewidth]{./figures_experiments/reports/attn_analysis_expr_first/attn_8b_one_hop_expr_first/query_region/attn_query_region_scores.png}
        \caption{Query (One-hop)}
    \end{subfigure}

    \vspace{0.2em}

    \begin{subfigure}[b]{0.285\textwidth}
        \centering
        \includegraphics[width=\linewidth]{./figures_experiments/reports/attn_analysis_expr_first/attn_8b_two_hop_expr_first/facts_region/attn_facts_region_scores.png}
        \caption{Facts (Two-hop)}
    \end{subfigure}
    \hfill
    \begin{subfigure}[b]{0.285\textwidth}
        \centering
        \includegraphics[width=\linewidth]{./figures_experiments/reports/attn_analysis_expr_first/attn_8b_two_hop_expr_first/expression_region/attn_expression_region_scores.png}
        \caption{Expression (Two-hop)}
    \end{subfigure}
    \hfill
    \begin{subfigure}[b]{0.285\textwidth}
        \centering
        \includegraphics[width=\linewidth]{./figures_experiments/reports/attn_analysis_expr_first/attn_8b_two_hop_expr_first/query_region/attn_query_region_scores.png}
        \caption{Query (Two-hop)}
    \end{subfigure}
    \caption{\textbf{Complete \textit{Qwen3-8B} attention-region gallery under \texttt{expr\_first}.} Relative to the MLP branch, attention patching shifts more mass toward middle-layer routing. The Facts Region is no longer purely early-layer dominated in one-hop, while the Query Region broadens into a larger middle-to-late control surface. This supports the interpretation that attention primarily redistributes and routes information once prompt order front-loads the expression.}
    \label{fig:qwen3_8b_attn_expr_gallery}
\end{figure*}

\begin{figure*}[t]
    \centering
    \begin{subfigure}[b]{0.285\textwidth}
        \centering
        \includegraphics[width=\linewidth]{./figures_experiments/reports/attn_analysis_expr_first/attn_14b_one_hop_expr_first/facts_region/attn_facts_region_scores.png}
        \caption{Facts (One-hop)}
    \end{subfigure}
    \hfill
    \begin{subfigure}[b]{0.285\textwidth}
        \centering
        \includegraphics[width=\linewidth]{./figures_experiments/reports/attn_analysis_expr_first/attn_14b_one_hop_expr_first/expression_region/attn_expression_region_scores.png}
        \caption{Expression (One-hop)}
    \end{subfigure}
    \hfill
    \begin{subfigure}[b]{0.285\textwidth}
        \centering
        \includegraphics[width=\linewidth]{./figures_experiments/reports/attn_analysis_expr_first/attn_14b_one_hop_expr_first/query_region/attn_query_region_scores.png}
        \caption{Query (One-hop)}
    \end{subfigure}

    \vspace{0.2em}

    \begin{subfigure}[b]{0.285\textwidth}
        \centering
        \includegraphics[width=\linewidth]{./figures_experiments/reports/attn_analysis_expr_first/attn_14b_two_hop_expr_first/facts_region/attn_facts_region_scores.png}
        \caption{Facts (Two-hop)}
    \end{subfigure}
    \hfill
    \begin{subfigure}[b]{0.285\textwidth}
        \centering
        \includegraphics[width=\linewidth]{./figures_experiments/reports/attn_analysis_expr_first/attn_14b_two_hop_expr_first/expression_region/attn_expression_region_scores.png}
        \caption{Expression (Two-hop)}
    \end{subfigure}
    \hfill
    \begin{subfigure}[b]{0.285\textwidth}
        \centering
        \includegraphics[width=\linewidth]{./figures_experiments/reports/attn_analysis_expr_first/attn_14b_two_hop_expr_first/query_region/attn_query_region_scores.png}
        \caption{Query (Two-hop)}
    \end{subfigure}
    \caption{\textbf{Complete \textit{Qwen3-14B} attention-region gallery under \texttt{expr\_first}.} The larger attention branch shows the same structure as 8B, but with an even clearer mid-layer Facts bulge in one-hop and a stronger middle-band Query component in two-hop. These curves make explicit that attention is the branch most sensitive to prompt-layout changes, especially in how it redistributes factual information away from the earliest layers.}
    \label{fig:qwen3_14b_attn_expr_gallery}
\end{figure*}

\begin{figure*}[t]
    \centering
    \begin{subfigure}[b]{0.285\textwidth}
        \centering
        \includegraphics[width=\linewidth]{./figures_experiments/reports/attn_mlp_analysis_expr_first/attn_mlp_8b_one_hop_expr_first/facts_region/attn_mlp_facts_region_scores.png}
        \caption{Facts (One-hop)}
    \end{subfigure}
    \hfill
    \begin{subfigure}[b]{0.285\textwidth}
        \centering
        \includegraphics[width=\linewidth]{./figures_experiments/reports/attn_mlp_analysis_expr_first/attn_mlp_8b_one_hop_expr_first/expression_region/attn_mlp_expression_region_scores.png}
        \caption{Expression (One-hop)}
    \end{subfigure}
    \hfill
    \begin{subfigure}[b]{0.285\textwidth}
        \centering
        \includegraphics[width=\linewidth]{./figures_experiments/reports/attn_mlp_analysis_expr_first/attn_mlp_8b_one_hop_expr_first/query_region/attn_mlp_query_region_scores.png}
        \caption{Query (One-hop)}
    \end{subfigure}

    \vspace{0.2em}

    \begin{subfigure}[b]{0.285\textwidth}
        \centering
        \includegraphics[width=\linewidth]{./figures_experiments/reports/attn_mlp_analysis_expr_first/attn_mlp_8b_two_hop_expr_first/facts_region/attn_mlp_facts_region_scores.png}
        \caption{Facts (Two-hop)}
    \end{subfigure}
    \hfill
    \begin{subfigure}[b]{0.285\textwidth}
        \centering
        \includegraphics[width=\linewidth]{./figures_experiments/reports/attn_mlp_analysis_expr_first/attn_mlp_8b_two_hop_expr_first/expression_region/attn_mlp_expression_region_scores.png}
        \caption{Expression (Two-hop)}
    \end{subfigure}
    \hfill
    \begin{subfigure}[b]{0.285\textwidth}
        \centering
        \includegraphics[width=\linewidth]{./figures_experiments/reports/attn_mlp_analysis_expr_first/attn_mlp_8b_two_hop_expr_first/query_region/attn_mlp_query_region_scores.png}
        \caption{Query (Two-hop)}
    \end{subfigure}
    \caption{\textbf{Complete \textit{Qwen3-8B} joint attention+MLP gallery under \texttt{expr\_first}.} The joint intervention remains more stage-like than the pure attention branch. Facts and expressions both retain strong early peaks, while Query continues to broaden over the middle and late bands. The figure therefore shows that coupling routing with nonlinear transformation partially stabilizes the staged decomposition under prompt reordering.}
    \label{fig:qwen3_8b_attn_mlp_expr_gallery}
\end{figure*}

\begin{figure*}[t]
    \centering
    \begin{subfigure}[b]{0.285\textwidth}
        \centering
        \includegraphics[width=\linewidth]{./figures_experiments/reports/attn_mlp_analysis_expr_first/attn_mlp_14b_one_hop_expr_first/facts_region/attn_mlp_facts_region_scores.png}
        \caption{Facts (One-hop)}
    \end{subfigure}
    \hfill
    \begin{subfigure}[b]{0.285\textwidth}
        \centering
        \includegraphics[width=\linewidth]{./figures_experiments/reports/attn_mlp_analysis_expr_first/attn_mlp_14b_one_hop_expr_first/expression_region/attn_mlp_expression_region_scores.png}
        \caption{Expression (One-hop)}
    \end{subfigure}
    \hfill
    \begin{subfigure}[b]{0.285\textwidth}
        \centering
        \includegraphics[width=\linewidth]{./figures_experiments/reports/attn_mlp_analysis_expr_first/attn_mlp_14b_one_hop_expr_first/query_region/attn_mlp_query_region_scores.png}
        \caption{Query (One-hop)}
    \end{subfigure}

    \vspace{0.2em}

    \begin{subfigure}[b]{0.285\textwidth}
        \centering
        \includegraphics[width=\linewidth]{./figures_experiments/reports/attn_mlp_analysis_expr_first/attn_mlp_14b_two_hop_expr_first/facts_region/attn_mlp_facts_region_scores.png}
        \caption{Facts (Two-hop)}
    \end{subfigure}
    \hfill
    \begin{subfigure}[b]{0.285\textwidth}
        \centering
        \includegraphics[width=\linewidth]{./figures_experiments/reports/attn_mlp_analysis_expr_first/attn_mlp_14b_two_hop_expr_first/expression_region/attn_mlp_expression_region_scores.png}
        \caption{Expression (Two-hop)}
    \end{subfigure}
    \hfill
    \begin{subfigure}[b]{0.285\textwidth}
        \centering
        \includegraphics[width=\linewidth]{./figures_experiments/reports/attn_mlp_analysis_expr_first/attn_mlp_14b_two_hop_expr_first/query_region/attn_mlp_query_region_scores.png}
        \caption{Query (Two-hop)}
    \end{subfigure}
    \caption{\textbf{Complete \textit{Qwen3-14B} joint attention+MLP gallery under \texttt{expr\_first}.} The 14B joint branch amplifies the one-hop Expression peak while keeping Facts strongly early-layer concentrated. Two-hop curves remain broader, but still preserve the same order: Facts and Expression dominate the front half of the network, and Query remains the principal sink for later integration.}
    \label{fig:qwen3_14b_attn_mlp_expr_gallery}
\end{figure*}

%\begin{figure*}[p]
%    \centering
%    \begin{subfigure}[b]{0.29\textwidth}
%        \centering
%        \includegraphics[width=\linewidth]{./figures_experiments/reports_llama3/mlp_analysis_facts_first/comparison/bcr_stacked.png}
%        \caption{Llama: MLP}
%    \end{subfigure}
%    \hfill
%    \begin{subfigure}[b]{0.29\textwidth}
%        \centering
%        \includegraphics[width=\linewidth]{./figures_experiments/reports_llama3/attn_analysis_facts_first/comparison/bcr_stacked.png}
%        \caption{Llama: Attention}
%    \end{subfigure}
%    \hfill
%    \begin{subfigure}[b]{0.29\textwidth}
%        \centering
%        \includegraphics[width=\linewidth]{./figures_experiments/reports_llama3/attn_mlp_analysis_facts_first/comparison/bcr_stacked.png}
%        \caption{Llama: Attention+MLP}
%    \end{subfigure}
%
%    \vspace{0.2em}
%
%    \begin{subfigure}[b]{0.29\textwidth}
%        \centering
%        \includegraphics[width=\linewidth]{./figures_experiments/reports_mistral/mlp_analysis_facts_first/comparison/bcr_stacked.png}
%        \caption{Mistral: MLP}
%    \end{subfigure}
%    \hfill
%    \begin{subfigure}[b]{0.29\textwidth}
%        \centering
%        \includegraphics[width=\linewidth]{./figures_experiments/reports_mistral/attn_analysis_facts_first/comparison/bcr_stacked.png}
%        \caption{Mistral: Attention}
%    \end{subfigure}
%    \hfill
%    \begin{subfigure}[b]{0.29\textwidth}
%        \centering
%        \includegraphics[width=\linewidth]{./figures_experiments/reports_mistral/attn_mlp_analysis_facts_first/comparison/bcr_stacked.png}
%        \caption{Mistral: Attention+MLP}
%    \end{subfigure}
%    \caption{\textbf{Cross-model stacked BCR summaries under \texttt{facts\_first}.} Both Llama and Mistral preserve the same coarse facts-to-expression-to-query ordering, but the stacked decompositions reveal much heavier query-facing mass than in Qwen3. Llama spreads more mass into the Query Region across all three branches, while Mistral retains a comparatively strong middle-band Query component even in the MLP branch.}
%    \label{fig:cross_model_bcr_stacked_facts}
%\end{figure*}

%\begin{figure*}[p]
%    \centering
%    \begin{subfigure}[b]{0.29\textwidth}
%        \centering
%        \includegraphics[width=\linewidth]{./figures_experiments/reports_llama3/mlp_analysis_expr_first/comparison/bcr_stacked.png}
%        \caption{Llama: MLP}
%    \end{subfigure}
%    \hfill
%    \begin{subfigure}[b]{0.29\textwidth}
%        \centering
%        \includegraphics[width=\linewidth]{./figures_experiments/reports_llama3/attn_analysis_expr_first/comparison/bcr_stacked.png}
%        \caption{Llama: Attention}
%    \end{subfigure}
%    \hfill
%    \begin{subfigure}[b]{0.29\textwidth}
%        \centering
%        \includegraphics[width=\linewidth]{./figures_experiments/reports_llama3/attn_mlp_analysis_expr_first/comparison/bcr_stacked.png}
%        \caption{Llama: Attention+MLP}
%    \end{subfigure}
%
%    \vspace{0.2em}
%
%    \begin{subfigure}[b]{0.29\textwidth}
%        \centering
%        \includegraphics[width=\linewidth]{./figures_experiments/reports_mistral/mlp_analysis_expr_first/comparison/bcr_stacked.png}
%        \caption{Mistral: MLP}
%    \end{subfigure}
%    \hfill
%    \begin{subfigure}[b]{0.29\textwidth}
%        \centering
%        \includegraphics[width=\linewidth]{./figures_experiments/reports_mistral/attn_analysis_expr_first/comparison/bcr_stacked.png}
%        \caption{Mistral: Attention}
%    \end{subfigure}
%    \hfill
%    \begin{subfigure}[b]{0.29\textwidth}
%        \centering
%        \includegraphics[width=\linewidth]{./figures_experiments/reports_mistral/attn_mlp_analysis_expr_first/comparison/bcr_stacked.png}
%        \caption{Mistral: Attention+MLP}
%    \end{subfigure}
%    \caption{\textbf{Cross-model stacked BCR summaries under \texttt{expr\_first}.} Prompt reordering again increases early Expression mass, but the cross-model gap persists. Llama keeps a notably large early-plus-late Query component, while Mistral preserves a broad Query distribution across the back half of the network. These summaries reinforce the interpretation that both non-Qwen families rely more heavily on query-facing scaffolding during final integration.}
%    \label{fig:cross_model_bcr_stacked_expr}
%\end{figure*}

\paragraph{How to compare the cross-model region figures.}
In the cross-model panels below, the safest comparison is to hold the branch fixed and then compare families: first compare Qwen3/Llama/Mistral within MLP, then within attention, then within the joint branch. A family with more late Query mass is relying more heavily on query-facing scaffolding during final integration, whereas a family with cleaner early Facts and middle Expression concentration is implementing a sharper staged decomposition.

The full cross-model MLP galleries sharpen the same qualitative point. In Llama, even the core MLP branch allocates unusually large mass to the Query Region: under one-hop \texttt{facts\_first}, the query profile is split across early ($0.51$), middle ($0.27$), and late ($0.22$) bands rather than being primarily a late aggregator. Mistral looks somewhat cleaner than Llama, but still remains more query-heavy than Qwen3, with one-hop \texttt{facts\_first} query shares of $0.48$ / $0.37$ / $0.15$ (early / middle / late). The full curves below make these differences visible at layer resolution rather than only through compact band summaries.

\begin{figure*}[t]
    \centering
    \begin{subfigure}[b]{0.285\textwidth}
        \centering
        \includegraphics[width=\linewidth]{./figures_experiments/reports_llama3/mlp_analysis_facts_first/mlp_one_hop_facts_first/facts_region/mlp_facts_region_scores.png}
        \caption{Facts (One-hop)}
    \end{subfigure}
    \hfill
    \begin{subfigure}[b]{0.285\textwidth}
        \centering
        \includegraphics[width=\linewidth]{./figures_experiments/reports_llama3/mlp_analysis_facts_first/mlp_one_hop_facts_first/expression_region/mlp_expression_region_scores.png}
        \caption{Expression (One-hop)}
    \end{subfigure}
    \hfill
    \begin{subfigure}[b]{0.285\textwidth}
        \centering
        \includegraphics[width=\linewidth]{./figures_experiments/reports_llama3/mlp_analysis_facts_first/mlp_one_hop_facts_first/query_region/mlp_query_region_scores.png}
        \caption{Query (One-hop)}
    \end{subfigure}

    \vspace{0.2em}

    \begin{subfigure}[b]{0.285\textwidth}
        \centering
        \includegraphics[width=\linewidth]{./figures_experiments/reports_llama3/mlp_analysis_facts_first/mlp_two_hop_facts_first/facts_region/mlp_facts_region_scores.png}
        \caption{Facts (Two-hop)}
    \end{subfigure}
    \hfill
    \begin{subfigure}[b]{0.285\textwidth}
        \centering
        \includegraphics[width=\linewidth]{./figures_experiments/reports_llama3/mlp_analysis_facts_first/mlp_two_hop_facts_first/expression_region/mlp_expression_region_scores.png}
        \caption{Expression (Two-hop)}
    \end{subfigure}
    \hfill
    \begin{subfigure}[b]{0.285\textwidth}
        \centering
        \includegraphics[width=\linewidth]{./figures_experiments/reports_llama3/mlp_analysis_facts_first/mlp_two_hop_facts_first/query_region/mlp_query_region_scores.png}
        \caption{Query (Two-hop)}
    \end{subfigure}
    \caption{\textbf{Complete Llama MLP region gallery under \texttt{facts\_first}.} Llama preserves the broad staged template but with a much heavier Query Region than Qwen3. The query curve is already substantial in one-hop middle layers and remains comparably strong in late layers, which is consistent with a more output-position-centered computation strategy.}
    \label{fig:llama_mlp_facts_gallery}
\end{figure*}

\begin{figure*}[t]
    \centering
    \begin{subfigure}[b]{0.285\textwidth}
        \centering
        \includegraphics[width=\linewidth]{./figures_experiments/reports_llama3/mlp_analysis_expr_first/mlp_one_hop_expr_first/facts_region/mlp_facts_region_scores.png}
        \caption{Facts (One-hop)}
    \end{subfigure}
    \hfill
    \begin{subfigure}[b]{0.285\textwidth}
        \centering
        \includegraphics[width=\linewidth]{./figures_experiments/reports_llama3/mlp_analysis_expr_first/mlp_one_hop_expr_first/expression_region/mlp_expression_region_scores.png}
        \caption{Expression (One-hop)}
    \end{subfigure}
    \hfill
    \begin{subfigure}[b]{0.285\textwidth}
        \centering
        \includegraphics[width=\linewidth]{./figures_experiments/reports_llama3/mlp_analysis_expr_first/mlp_one_hop_expr_first/query_region/mlp_query_region_scores.png}
        \caption{Query (One-hop)}
    \end{subfigure}

    \vspace{0.2em}

    \begin{subfigure}[b]{0.285\textwidth}
        \centering
        \includegraphics[width=\linewidth]{./figures_experiments/reports_llama3/mlp_analysis_expr_first/mlp_two_hop_expr_first/facts_region/mlp_facts_region_scores.png}
        \caption{Facts (Two-hop)}
    \end{subfigure}
    \hfill
    \begin{subfigure}[b]{0.285\textwidth}
        \centering
        \includegraphics[width=\linewidth]{./figures_experiments/reports_llama3/mlp_analysis_expr_first/mlp_two_hop_expr_first/expression_region/mlp_expression_region_scores.png}
        \caption{Expression (Two-hop)}
    \end{subfigure}
    \hfill
    \begin{subfigure}[b]{0.285\textwidth}
        \centering
        \includegraphics[width=\linewidth]{./figures_experiments/reports_llama3/mlp_analysis_expr_first/mlp_two_hop_expr_first/query_region/mlp_query_region_scores.png}
        \caption{Query (Two-hop)}
    \end{subfigure}
    \caption{\textbf{Complete Llama MLP region gallery under \texttt{expr\_first}.} Prompt reordering shifts expression processing earlier as expected, but the Query Region remains unusually broad and strong. Even after moving the queried expression to the front, Llama still allocates a large amount of MLP sensitivity to the query-facing suffix and answer anchor.}
    \label{fig:llama_mlp_expr_gallery}
\end{figure*}

\begin{figure*}[t]
    \centering
    \begin{subfigure}[b]{0.285\textwidth}
        \centering
        \includegraphics[width=\linewidth]{./figures_experiments/reports_mistral/mlp_analysis_facts_first/mlp_one_hop_facts_first/facts_region/mlp_facts_region_scores.png}
        \caption{Facts (One-hop)}
    \end{subfigure}
    \hfill
    \begin{subfigure}[b]{0.285\textwidth}
        \centering
        \includegraphics[width=\linewidth]{./figures_experiments/reports_mistral/mlp_analysis_facts_first/mlp_one_hop_facts_first/expression_region/mlp_expression_region_scores.png}
        \caption{Expression (One-hop)}
    \end{subfigure}
    \hfill
    \begin{subfigure}[b]{0.285\textwidth}
        \centering
        \includegraphics[width=\linewidth]{./figures_experiments/reports_mistral/mlp_analysis_facts_first/mlp_one_hop_facts_first/query_region/mlp_query_region_scores.png}
        \caption{Query (One-hop)}
    \end{subfigure}

    \vspace{0.2em}

    \begin{subfigure}[b]{0.285\textwidth}
        \centering
        \includegraphics[width=\linewidth]{./figures_experiments/reports_mistral/mlp_analysis_facts_first/mlp_two_hop_facts_first/facts_region/mlp_facts_region_scores.png}
        \caption{Facts (Two-hop)}
    \end{subfigure}
    \hfill
    \begin{subfigure}[b]{0.285\textwidth}
        \centering
        \includegraphics[width=\linewidth]{./figures_experiments/reports_mistral/mlp_analysis_facts_first/mlp_two_hop_facts_first/expression_region/mlp_expression_region_scores.png}
        \caption{Expression (Two-hop)}
    \end{subfigure}
    \hfill
    \begin{subfigure}[b]{0.285\textwidth}
        \centering
        \includegraphics[width=\linewidth]{./figures_experiments/reports_mistral/mlp_analysis_facts_first/mlp_two_hop_facts_first/query_region/mlp_query_region_scores.png}
        \caption{Query (Two-hop)}
    \end{subfigure}
    \caption{\textbf{Complete Mistral MLP region gallery under \texttt{facts\_first}.} Mistral retains a clearer early-Facts and early-Expression ordering than Llama, but again assigns more mass to the Query Region than Qwen3. The one-hop Query curve remains substantial throughout the back half of the network, and the two-hop curves further broaden this integration profile.}
    \label{fig:mistral_mlp_facts_gallery}
\end{figure*}

\begin{figure*}[t]
    \centering
    \begin{subfigure}[b]{0.285\textwidth}
        \centering
        \includegraphics[width=\linewidth]{./figures_experiments/reports_mistral/mlp_analysis_expr_first/mlp_one_hop_expr_first/facts_region/mlp_facts_region_scores.png}
        \caption{Facts (One-hop)}
    \end{subfigure}
    \hfill
    \begin{subfigure}[b]{0.285\textwidth}
        \centering
        \includegraphics[width=\linewidth]{./figures_experiments/reports_mistral/mlp_analysis_expr_first/mlp_one_hop_expr_first/expression_region/mlp_expression_region_scores.png}
        \caption{Expression (One-hop)}
    \end{subfigure}
    \hfill
    \begin{subfigure}[b]{0.285\textwidth}
        \centering
        \includegraphics[width=\linewidth]{./figures_experiments/reports_mistral/mlp_analysis_expr_first/mlp_one_hop_expr_first/query_region/mlp_query_region_scores.png}
        \caption{Query (One-hop)}
    \end{subfigure}

    \vspace{0.2em}

    \begin{subfigure}[b]{0.285\textwidth}
        \centering
        \includegraphics[width=\linewidth]{./figures_experiments/reports_mistral/mlp_analysis_expr_first/mlp_two_hop_expr_first/facts_region/mlp_facts_region_scores.png}
        \caption{Facts (Two-hop)}
    \end{subfigure}
    \hfill
    \begin{subfigure}[b]{0.285\textwidth}
        \centering
        \includegraphics[width=\linewidth]{./figures_experiments/reports_mistral/mlp_analysis_expr_first/mlp_two_hop_expr_first/expression_region/mlp_expression_region_scores.png}
        \caption{Expression (Two-hop)}
    \end{subfigure}
    \hfill
    \begin{subfigure}[b]{0.285\textwidth}
        \centering
        \includegraphics[width=\linewidth]{./figures_experiments/reports_mistral/mlp_analysis_expr_first/mlp_two_hop_expr_first/query_region/mlp_query_region_scores.png}
        \caption{Query (Two-hop)}
    \end{subfigure}
    \caption{\textbf{Complete Mistral MLP region gallery under \texttt{expr\_first}.} Mistral again shifts expression processing forward under prompt reordering, but the Query Region remains broad and relatively heavy. Compared with Qwen3, the figure suggests a weaker separation between semantic content processing and final answer integration.}
    \label{fig:mistral_mlp_expr_gallery}
\end{figure*}

\subsection{Additional Token Diagnostics}
\label{sec:appendix_extra_token_gallery}

The token diagnostics below sharpen two points left implicit in the main appendix. First, the refined mean and signed plots separate large-magnitude categories from categories that consistently move the model toward the correct answer. Second, the cross-model contrast remains stable across different normalizations. In unsigned totals, Qwen3 allocates substantially more mass to fact tokens than either non-Qwen family, whereas Llama and Mistral put relatively more mass on \texttt{query\_token} (under \texttt{facts\_first}: \textit{Qwen3-8B} \texttt{facts\_value} $2.47$ vs. \texttt{query\_token} $2.22$; Llama $0.72$ vs. $1.63$; Mistral $0.75$ vs. $1.04$). In signed totals, Qwen3 remains fact-dominant, Llama remains query-dominant, and Mistral falls between those extremes: under signed \texttt{facts\_first}, its \texttt{query\_token} and \texttt{facts\_value} totals are nearly identical ($0.731$ vs. $0.729$).

\paragraph{How to read these token figures.}
The refined \textbf{sum} plots answer a total-contribution question and therefore reward categories that accumulate substantial mass over multiple tokens, whereas the refined \textbf{mean} plots normalize by category size and answer a per-token-importance question. Signed variants preserve whether the intervention shifts the model in the correct direction, so a category can be large in unsigned magnitude but less dominant once directional consistency is required.

For Qwen3, the retained refined views are given in Figures~\ref{fig:qwen_token_refined_mean}, \ref{fig:qwen_token_refined_sum_signed}, and \ref{fig:qwen_token_refined_mean_signed}. The corresponding cross-model comparisons are collected in Figures~\ref{fig:cross_model_token_refined_mean}, \ref{fig:cross_model_token_refined_sum_signed}, and \ref{fig:cross_model_token_refined_mean_signed}.

%\begin{figure*}[p]
%    \centering
%    \begin{subfigure}[b]{0.44\textwidth}
%        \centering
%        \includegraphics[width=\linewidth]{./figures_experiments/reports/token_analysis/8b_facts_first_all_hop_raw/heatmap_simple.png}
%        \caption{\textit{Qwen3-8B}, \texttt{facts\_first}}
%    \end{subfigure}
%    \hfill
%    \begin{subfigure}[b]{0.44\textwidth}
%        \centering
%        \includegraphics[width=\linewidth]{./figures_experiments/reports/token_analysis/14b_facts_first_all_hop_raw/heatmap_simple.png}
%        \caption{\textit{Qwen3-14B}, \texttt{facts\_first}}
%    \end{subfigure}
%
%    \vspace{0.2em}
%
%    \begin{subfigure}[b]{0.44\textwidth}
%        \centering
%        \includegraphics[width=\linewidth]{./figures_experiments/reports/token_analysis/8b_expr_first_all_hop_raw/heatmap_simple.png}
%        \caption{\textit{Qwen3-8B}, \texttt{expr\_first}}
%    \end{subfigure}
%    \hfill
%    \begin{subfigure}[b]{0.44\textwidth}
%        \centering
%        \includegraphics[width=\linewidth]{./figures_experiments/reports/token_analysis/14b_expr_first_all_hop_raw/heatmap_simple.png}
%        \caption{\textit{Qwen3-14B}, \texttt{expr\_first}}
%    \end{subfigure}
%    \caption{\textbf{Qwen3 raw token heatmaps across scales and prompt orders.} The raw category-by-layer heatmaps make the same stage structure visible before refinement. \texttt{facts\_first} concentrates large early-layer effects on fact tokens, while \texttt{expr\_first} shifts more structure toward query-facing and derived-assignment categories. The 14B model keeps the same topology while producing a slightly cleaner early-to-late separation.}
%    \label{fig:qwen_token_raw_heatmaps}
%\end{figure*}

\begin{figure*}[t]
    \centering
    \begin{subfigure}[b]{0.44\textwidth}
        \centering
        \includegraphics[width=\linewidth]{./figures_experiments/reports/token_analysis/8b_facts_first_all_hop_refined/refined_by_stage_mean.png}
        \caption{\textit{Qwen3-8B}, \texttt{facts\_first}}
    \end{subfigure}
    \hfill
    \begin{subfigure}[b]{0.44\textwidth}
        \centering
        \includegraphics[width=\linewidth]{./figures_experiments/reports/token_analysis/14b_facts_first_all_hop_refined/refined_by_stage_mean.png}
        \caption{\textit{Qwen3-14B}, \texttt{facts\_first}}
    \end{subfigure}

    \vspace{0.2em}

    \begin{subfigure}[b]{0.44\textwidth}
        \centering
        \includegraphics[width=\linewidth]{./figures_experiments/reports/token_analysis/8b_expr_first_all_hop_refined/refined_by_stage_mean.png}
        \caption{\textit{Qwen3-8B}, \texttt{expr\_first}}
    \end{subfigure}
    \hfill
    \begin{subfigure}[b]{0.44\textwidth}
        \centering
        \includegraphics[width=\linewidth]{./figures_experiments/reports/token_analysis/14b_expr_first_all_hop_refined/refined_by_stage_mean.png}
        \caption{\textit{Qwen3-14B}, \texttt{expr\_first}}
    \end{subfigure}
    \caption{\textbf{Qwen3 refined mean token diagnostics.} Averaging by token rather than summing by category preserves the same hierarchy as the sum-based plots while shrinking the advantage of categories that occupy more positions. Facts and query remain the dominant token types in all four settings, which shows that the main result is not an artifact of category cardinality.}
    \label{fig:qwen_token_refined_mean}
\end{figure*}

\begin{figure*}[t]
    \centering
    \begin{subfigure}[b]{0.44\textwidth}
        \centering
        \includegraphics[width=\linewidth]{./figures_experiments/reports/token_analysis/8b_facts_first_all_hop_refined/refined_by_stage_sum_signed.png}
        \caption{\textit{Qwen3-8B}, \texttt{facts\_first}}
    \end{subfigure}
    \hfill
    \begin{subfigure}[b]{0.44\textwidth}
        \centering
        \includegraphics[width=\linewidth]{./figures_experiments/reports/token_analysis/14b_facts_first_all_hop_refined/refined_by_stage_sum_signed.png}
        \caption{\textit{Qwen3-14B}, \texttt{facts\_first}}
    \end{subfigure}

    \vspace{0.2em}

    \begin{subfigure}[b]{0.44\textwidth}
        \centering
        \includegraphics[width=\linewidth]{./figures_experiments/reports/token_analysis/8b_expr_first_all_hop_refined/refined_by_stage_sum_signed.png}
        \caption{\textit{Qwen3-8B}, \texttt{expr\_first}}
    \end{subfigure}
    \hfill
    \begin{subfigure}[b]{0.44\textwidth}
        \centering
        \includegraphics[width=\linewidth]{./figures_experiments/reports/token_analysis/14b_expr_first_all_hop_refined/refined_by_stage_sum_signed.png}
        \caption{\textit{Qwen3-14B}, \texttt{expr\_first}}
    \end{subfigure}
    \caption{\textbf{Qwen3 refined signed-sum token diagnostics.} These plots retain the sign of the token-level shift rather than using absolute magnitude. The dominant categories remain the same as in the unsigned plots, which shows that fact and query tokens are not merely high-variance sites: they push the model in a directionally meaningful way for the final decision.}
    \label{fig:qwen_token_refined_sum_signed}
\end{figure*}

\begin{figure*}[t]
    \centering
    \begin{subfigure}[b]{0.44\textwidth}
        \centering
        \includegraphics[width=\linewidth]{./figures_experiments/reports/token_analysis/8b_facts_first_all_hop_refined/refined_by_stage_mean_signed.png}
        \caption{\textit{Qwen3-8B}, \texttt{facts\_first}}
    \end{subfigure}
    \hfill
    \begin{subfigure}[b]{0.44\textwidth}
        \centering
        \includegraphics[width=\linewidth]{./figures_experiments/reports/token_analysis/14b_facts_first_all_hop_refined/refined_by_stage_mean_signed.png}
        \caption{\textit{Qwen3-14B}, \texttt{facts\_first}}
    \end{subfigure}

    \vspace{0.2em}

    \begin{subfigure}[b]{0.44\textwidth}
        \centering
        \includegraphics[width=\linewidth]{./figures_experiments/reports/token_analysis/8b_expr_first_all_hop_refined/refined_by_stage_mean_signed.png}
        \caption{\textit{Qwen3-8B}, \texttt{expr\_first}}
    \end{subfigure}
    \hfill
    \begin{subfigure}[b]{0.44\textwidth}
        \centering
        \includegraphics[width=\linewidth]{./figures_experiments/reports/token_analysis/14b_expr_first_all_hop_refined/refined_by_stage_mean_signed.png}
        \caption{\textit{Qwen3-14B}, \texttt{expr\_first}}
    \end{subfigure}
    \caption{\textbf{Qwen3 refined signed-mean token diagnostics.} The signed mean-per-token view confirms that the facts-versus-query balance is robust to both normalization choices. Under \texttt{expr\_first}, the facts category remains directionally dominant in 14B even though query-facing mass is more visible in the unsigned plots, indicating that prompt reordering broadens computation without eliminating fact access.}
    \label{fig:qwen_token_refined_mean_signed}
\end{figure*}

%\begin{figure*}[p]
%    \centering
%    \begin{subfigure}[b]{0.44\textwidth}
%        \centering
%        \includegraphics[width=\linewidth]{./figures_experiments/reports_llama3/token_analysis_facts_first_raw/heatmap_simple.png}
%        \caption{\textit{Llama-3.1-8B-Instruct}, \texttt{facts\_first}}
%    \end{subfigure}
%    \hfill
%    \begin{subfigure}[b]{0.44\textwidth}
%        \centering
%        \includegraphics[width=\linewidth]{./figures_experiments/reports_llama3/token_analysis_expr_first_raw/heatmap_simple.png}
%        \caption{\textit{Llama-3.1-8B-Instruct}, \texttt{expr\_first}}
%    \end{subfigure}
%
%    \vspace{0.2em}
%
%    \begin{subfigure}[b]{0.44\textwidth}
%        \centering
%        \includegraphics[width=\linewidth]{./figures_experiments/reports_mistral/token_analysis_facts_first_raw/heatmap_simple.png}
%        \caption{\textit{Mistral-7B-Instruct-v0.1}, \texttt{facts\_first}}
%    \end{subfigure}
%    \hfill
%    \begin{subfigure}[b]{0.44\textwidth}
%        \centering
%        \includegraphics[width=\linewidth]{./figures_experiments/reports_mistral/token_analysis_expr_first_raw/heatmap_simple.png}
%        \caption{\textit{Mistral-7B-Instruct-v0.1}, \texttt{expr\_first}}
%    \end{subfigure}
%    \caption{\textbf{Cross-model raw token heatmaps.} Before any category refinement, both non-Qwen families already show heavier query-facing structure than Qwen3. Llama is especially diffuse, with broad late-layer activity over the query-facing portion of the prompt, while Mistral keeps a somewhat cleaner early-to-late progression but still concentrates less sharply on fact tokens than Qwen3.}
%    \label{fig:cross_model_token_raw_heatmaps}
%\end{figure*}

\begin{figure*}[t]
    \centering
    \begin{subfigure}[b]{0.44\textwidth}
        \centering
        \includegraphics[width=\linewidth]{./figures_experiments/reports_llama3/token_analysis_facts_first_refined/refined_by_stage_mean.png}
        \caption{\textit{Llama-3.1-8B-Instruct}, \texttt{facts\_first}}
    \end{subfigure}
    \hfill
    \begin{subfigure}[b]{0.44\textwidth}
        \centering
        \includegraphics[width=\linewidth]{./figures_experiments/reports_llama3/token_analysis_expr_first_refined/refined_by_stage_mean.png}
        \caption{\textit{Llama-3.1-8B-Instruct}, \texttt{expr\_first}}
    \end{subfigure}

    \vspace{0.2em}

    \begin{subfigure}[b]{0.44\textwidth}
        \centering
        \includegraphics[width=\linewidth]{./figures_experiments/reports_mistral/token_analysis_facts_first_refined/refined_by_stage_mean.png}
        \caption{\textit{Mistral-7B-Instruct-v0.1}, \texttt{facts\_first}}
    \end{subfigure}
    \hfill
    \begin{subfigure}[b]{0.44\textwidth}
        \centering
        \includegraphics[width=\linewidth]{./figures_experiments/reports_mistral/token_analysis_expr_first_refined/refined_by_stage_mean.png}
        \caption{\textit{Mistral-7B-Instruct-v0.1}, \texttt{expr\_first}}
    \end{subfigure}
    \caption{\textbf{Cross-model refined mean token diagnostics.} In both non-Qwen families, \texttt{query\_token} remains the largest mean contribution under both prompt orders. This is a compact quantitative diagnostic of the more output-position-centered strategy already visible in the band-composition summaries.}
    \label{fig:cross_model_token_refined_mean}
\end{figure*}

\begin{figure*}[t]
    \centering
    \begin{subfigure}[b]{0.44\textwidth}
        \centering
        \includegraphics[width=\linewidth]{./figures_experiments/reports_llama3/token_analysis_facts_first_refined/refined_by_stage_sum_signed.png}
        \caption{\textit{Llama-3.1-8B-Instruct}, \texttt{facts\_first}}
    \end{subfigure}
    \hfill
    \begin{subfigure}[b]{0.44\textwidth}
        \centering
        \includegraphics[width=\linewidth]{./figures_experiments/reports_llama3/token_analysis_expr_first_refined/refined_by_stage_sum_signed.png}
        \caption{\textit{Llama-3.1-8B-Instruct}, \texttt{expr\_first}}
    \end{subfigure}

    \vspace{0.2em}

    \begin{subfigure}[b]{0.44\textwidth}
        \centering
        \includegraphics[width=\linewidth]{./figures_experiments/reports_mistral/token_analysis_facts_first_refined/refined_by_stage_sum_signed.png}
        \caption{\textit{Mistral-7B-Instruct-v0.1}, \texttt{facts\_first}}
    \end{subfigure}
    \hfill
    \begin{subfigure}[b]{0.44\textwidth}
        \centering
        \includegraphics[width=\linewidth]{./figures_experiments/reports_mistral/token_analysis_expr_first_refined/refined_by_stage_sum_signed.png}
        \caption{\textit{Mistral-7B-Instruct-v0.1}, \texttt{expr\_first}}
    \end{subfigure}
    \caption{\textbf{Cross-model refined signed-sum token diagnostics.} Retaining the sign sharpens the inter-family contrast, but not uniformly. Llama remains clearly query-dominant, Qwen3 remains clearly fact-dominant, and Mistral falls between the two: it is nearly balanced under \texttt{facts\_first} and becomes more query-heavy under \texttt{expr\_first}. This is why Mistral looks more query-facing than Qwen3 without being as one-sided as Llama in every signed view.}
    \label{fig:cross_model_token_refined_sum_signed}
\end{figure*}

\begin{figure*}[t]
    \centering
    \begin{subfigure}[b]{0.44\textwidth}
        \centering
        \includegraphics[width=\linewidth]{./figures_experiments/reports_llama3/token_analysis_facts_first_refined/refined_by_stage_mean_signed.png}
        \caption{\textit{Llama-3.1-8B-Instruct}, \texttt{facts\_first}}
    \end{subfigure}
    \hfill
    \begin{subfigure}[b]{0.44\textwidth}
        \centering
        \includegraphics[width=\linewidth]{./figures_experiments/reports_llama3/token_analysis_expr_first_refined/refined_by_stage_mean_signed.png}
        \caption{\textit{Llama-3.1-8B-Instruct}, \texttt{expr\_first}}
    \end{subfigure}

    \vspace{0.2em}

    \begin{subfigure}[b]{0.44\textwidth}
        \centering
        \includegraphics[width=\linewidth]{./figures_experiments/reports_mistral/token_analysis_facts_first_refined/refined_by_stage_mean_signed.png}
        \caption{\textit{Mistral-7B-Instruct-v0.1}, \texttt{facts\_first}}
    \end{subfigure}
    \hfill
    \begin{subfigure}[b]{0.44\textwidth}
        \centering
        \includegraphics[width=\linewidth]{./figures_experiments/reports_mistral/token_analysis_expr_first_refined/refined_by_stage_mean_signed.png}
        \caption{\textit{Mistral-7B-Instruct-v0.1}, \texttt{expr\_first}}
    \end{subfigure}
    \caption{\textbf{Cross-model refined signed-mean token diagnostics.} The per-token normalization leaves the same ordering intact: Llama and Mistral remain query-heavy even after accounting for category size. This confirms that the cross-model difference is not driven by having more query-region tokens in the prompt template, but by where the models actually place their causal sensitivity.}
    \label{fig:cross_model_token_refined_mean_signed}
\end{figure*}

\subsection{Additional Head Diagnostics}
\label{sec:appendix_extra_heads_gallery}

The omitted head-analysis figures extend the taxonomy line charts already shown in the appendix. The Step~1 discovery maps show whether candidate heads are spatially clustered before any role label is assigned. The Step~3 plots then expose both the magnitude and direction of the validation effect through absolute-ratio and raw dPD-shift curves. For Qwen3, these diagnostics align closely: candidate heads cluster in the same early/middle/late bands later recovered by the taxonomy, and role-targeted curves remain much more damaging than \textit{Random-Outside}. For Llama and Mistral, the same views show weaker separation and larger side effects from the random control.

\paragraph{How to read these head figures.}
The discovery maps should be read as \emph{candidate density over layers and heads}: clustering means that promising heads are already spatially organized before the taxonomy rules are applied. The absolute-ratio validation curves answer ``how damaging is this ablation relative to the model's original margin?'', so higher curves are worse. The raw dPD-shift curves add direction: more negative values mean the ablation pushes the model farther away from the correct answer, while values near zero indicate weak causal effect.

The Qwen3-specific head extensions appear in Figures~\ref{fig:qwen_heads_step1_gallery} and~\ref{fig:qwen_heads_dpd_shift_gallery}, while the corresponding cross-model diagnostics appear in Figures~\ref{fig:cross_model_heads_step1_gallery} and~\ref{fig:cross_model_heads_dpd_shift_gallery}.

\begin{figure*}[t]
    \centering
    \begin{subfigure}[b]{0.44\textwidth}
        \centering
        \includegraphics[width=\linewidth]{./figures_experiments/reports/heads_analysis_qwen3_8b_facts_first/classify/layer_head_role_distribution.png}
        \caption{\textit{Qwen3-8B}, \texttt{facts\_first}}
    \end{subfigure}
    \hfill
    \begin{subfigure}[b]{0.44\textwidth}
        \centering
        \includegraphics[width=\linewidth]{./figures_experiments/reports/heads_analysis_qwen3_8b_expr_first/classify/layer_head_role_distribution.png}
        \caption{\textit{Qwen3-8B}, \texttt{expr\_first}}
    \end{subfigure}

    \vspace{0.2em}

    \begin{subfigure}[b]{0.44\textwidth}
        \centering
        \includegraphics[width=\linewidth]{./figures_experiments/reports/heads_analysis_qwen3_14b_facts_first/classify/layer_head_role_distribution.png}
        \caption{\textit{Qwen3-14B}, \texttt{facts\_first}}
    \end{subfigure}
    \hfill
    \begin{subfigure}[b]{0.44\textwidth}
        \centering
        \includegraphics[width=\linewidth]{./figures_experiments/reports/heads_analysis_qwen3_14b_expr_first/classify/layer_head_role_distribution.png}
        \caption{\textit{Qwen3-14B}, \texttt{expr\_first}}
    \end{subfigure}
    \caption{\textbf{Qwen3 Step~1 discovery maps across scales and prompt orders.} Even before taxonomy assignment, the candidate heads are not uniformly spread across the network. Early-layer candidates cluster near boundary-sensitive regions, middle-layer candidates align with transmission bands, and later candidates shift toward fact-retrieval behavior. The \texttt{expr\_first} intervention increases the density of early-layer candidates, matching the higher abundance of Splitting Heads reported in Step~2.}
    \label{fig:qwen_heads_step1_gallery}
\end{figure*}

\begin{figure*}[t]
    \centering
    \begin{subfigure}[b]{0.44\textwidth}
        \centering
        \includegraphics[width=\linewidth]{./figures_experiments/reports/heads_analysis_qwen3_8b_facts_first/validation/plots/dpd_shift_curve.png}
        \caption{\textit{Qwen3-8B}, \texttt{facts\_first}}
    \end{subfigure}
    \hfill
    \begin{subfigure}[b]{0.44\textwidth}
        \centering
        \includegraphics[width=\linewidth]{./figures_experiments/reports/heads_analysis_qwen3_8b_expr_first/validation/plots/dpd_shift_curve.png}
        \caption{\textit{Qwen3-8B}, \texttt{expr\_first}}
    \end{subfigure}

    \vspace{0.2em}

    \begin{subfigure}[b]{0.44\textwidth}
        \centering
        \includegraphics[width=\linewidth]{./figures_experiments/reports/heads_analysis_qwen3_14b_facts_first/validation/plots/dpd_shift_curve.png}
        \caption{\textit{Qwen3-14B}, \texttt{facts\_first}}
    \end{subfigure}
    \hfill
    \begin{subfigure}[b]{0.44\textwidth}
        \centering
        \includegraphics[width=\linewidth]{./figures_experiments/reports/heads_analysis_qwen3_14b_expr_first/validation/plots/dpd_shift_curve.png}
        \caption{\textit{Qwen3-14B}, \texttt{expr\_first}}
    \end{subfigure}
    \caption{\textbf{Qwen3 Step~3 dPD-shift diagnostics.} The raw probability-difference shifts are again strongly negative for the discovered roles and remain close to zero for the random baseline. The effect is strongest in \textit{Qwen3-8B} \texttt{facts\_first}, where Fact-Retrieval heads cause the largest downward shift, but the same ordering remains visible in the other three settings.}
    \label{fig:qwen_heads_dpd_shift_gallery}
\end{figure*}

\begin{figure*}[t]
    \centering
    \begin{subfigure}[b]{0.44\textwidth}
        \centering
        \includegraphics[width=\linewidth]{./figures_experiments/reports_llama3/heads_analysis_facts_first/classify/layer_head_role_distribution.png}
        \caption{Llama, \texttt{facts\_first}}
    \end{subfigure}
    \hfill
    \begin{subfigure}[b]{0.44\textwidth}
        \centering
        \includegraphics[width=\linewidth]{./figures_experiments/reports_llama3/heads_analysis_expr_first/classify/layer_head_role_distribution.png}
        \caption{Llama, \texttt{expr\_first}}
    \end{subfigure}

    \vspace{0.2em}

    \begin{subfigure}[b]{0.44\textwidth}
        \centering
        \includegraphics[width=\linewidth]{./figures_experiments/reports_mistral/heads_analysis_facts_first/classify/layer_head_role_distribution.png}
        \caption{Mistral, \texttt{facts\_first}}
    \end{subfigure}
    \hfill
    \begin{subfigure}[b]{0.44\textwidth}
        \centering
        \includegraphics[width=\linewidth]{./figures_experiments/reports_mistral/heads_analysis_expr_first/classify/layer_head_role_distribution.png}
        \caption{Mistral, \texttt{expr\_first}}
    \end{subfigure}
    \caption{\textbf{Cross-model Step~1 discovery maps.} The same coarse early/middle/late clustering is still visible in Llama and Mistral, but the candidate sets are visibly more diffuse than in Qwen3. This foreshadows the weaker role separation later visible in the absolute-ratio and dPD-shift validation curves.}
    \label{fig:cross_model_heads_step1_gallery}
\end{figure*}

\begin{figure*}[t]
    \centering
    \begin{subfigure}[b]{0.44\textwidth}
        \centering
        \includegraphics[width=\linewidth]{./figures_experiments/reports_llama3/heads_analysis_facts_first/validation/plots/dpd_shift_curve.png}
        \caption{Llama, \texttt{facts\_first}}
    \end{subfigure}
    \hfill
    \begin{subfigure}[b]{0.44\textwidth}
        \centering
        \includegraphics[width=\linewidth]{./figures_experiments/reports_llama3/heads_analysis_expr_first/validation/plots/dpd_shift_curve.png}
        \caption{Llama, \texttt{expr\_first}}
    \end{subfigure}

    \vspace{0.2em}

    \begin{subfigure}[b]{0.44\textwidth}
        \centering
        \includegraphics[width=\linewidth]{./figures_experiments/reports_mistral/heads_analysis_facts_first/validation/plots/dpd_shift_curve.png}
        \caption{Mistral, \texttt{facts\_first}}
    \end{subfigure}
    \hfill
    \begin{subfigure}[b]{0.44\textwidth}
        \centering
        \includegraphics[width=\linewidth]{./figures_experiments/reports_mistral/heads_analysis_expr_first/validation/plots/dpd_shift_curve.png}
        \caption{Mistral, \texttt{expr\_first}}
    \end{subfigure}
    \caption{\textbf{Cross-model Step~3 dPD-shift diagnostics.} The raw dPD shifts confirm the same story. Llama's targeted interventions are harmful but shallow, and the random baseline can even move in the helpful direction. Mistral's targeted interventions remain sharply harmful, yet the random baseline also produces a noticeable positive displacement, consistent with a highly sensitive but less cleanly isolated head substrate.}
    \label{fig:cross_model_heads_dpd_shift_gallery}
\end{figure*}

\subsection{Further Region and Magnitude Atlases}
\label{sec:appendix_further_atlas}

This subsection collects the remaining retained region galleries. They are included to complete the omitted hop- and branch-specific views, especially where the earlier appendix sections only gave a partial picture.

\paragraph{Qwen3 facts-first supplements.}
The remaining Qwen3 figures fill three specific gaps: the omitted two-hop attention gallery, the omitted two-hop joint attention+MLP gallery, and the retained facts-first two-hop MLP gallery. Together they show that increased reasoning depth broadens the Query Region without changing the overall Facts $\rightarrow$ Expression $\rightarrow$ Query order.

\begin{figure*}[t]
    \centering
    \begin{subfigure}[b]{0.285\textwidth}
        \centering
        \includegraphics[width=\linewidth]{./figures_experiments/reports/attn_analysis/attn_8b_two_hop/facts_region/attn_facts_region_scores.png}
        \caption{\textit{Qwen3-8B}: Facts}
    \end{subfigure}
    \hfill
    \begin{subfigure}[b]{0.285\textwidth}
        \centering
        \includegraphics[width=\linewidth]{./figures_experiments/reports/attn_analysis/attn_8b_two_hop/expression_region/attn_expression_region_scores.png}
        \caption{\textit{Qwen3-8B}: Expression}
    \end{subfigure}
    \hfill
    \begin{subfigure}[b]{0.285\textwidth}
        \centering
        \includegraphics[width=\linewidth]{./figures_experiments/reports/attn_analysis/attn_8b_two_hop/query_region/attn_query_region_scores.png}
        \caption{\textit{Qwen3-8B}: Query}
    \end{subfigure}

    \vspace{0.2em}

    \begin{subfigure}[b]{0.285\textwidth}
        \centering
        \includegraphics[width=\linewidth]{./figures_experiments/reports/attn_analysis/attn_14b_two_hop/facts_region/attn_facts_region_scores.png}
        \caption{\textit{Qwen3-14B}: Facts}
    \end{subfigure}
    \hfill
    \begin{subfigure}[b]{0.285\textwidth}
        \centering
        \includegraphics[width=\linewidth]{./figures_experiments/reports/attn_analysis/attn_14b_two_hop/expression_region/attn_expression_region_scores.png}
        \caption{\textit{Qwen3-14B}: Expression}
    \end{subfigure}
    \hfill
    \begin{subfigure}[b]{0.285\textwidth}
        \centering
        \includegraphics[width=\linewidth]{./figures_experiments/reports/attn_analysis/attn_14b_two_hop/query_region/attn_query_region_scores.png}
        \caption{\textit{Qwen3-14B}: Query}
    \end{subfigure}
    \caption{\textbf{Qwen3 facts-first attention-region galleries for two-hop reasoning.} Two-hop attention patching preserves the same semantic order as one-hop, but widens the query-facing integration window. In both model sizes, the Facts Region remains early-dominant, the Expression Region splits between early and middle layers, and the Query Region acquires its largest mass in the middle-to-late bands.}
    \label{fig:qwen_attn_twohop_facts_gallery}
\end{figure*}

\begin{figure*}[t]
    \centering
    \begin{subfigure}[b]{0.285\textwidth}
        \centering
        \includegraphics[width=\linewidth]{./figures_experiments/reports/attn_mlp_analysis/attn_mlp_8b_two_hop/facts_region/attn_mlp_facts_region_scores.png}
        \caption{\textit{Qwen3-8B}: Facts}
    \end{subfigure}
    \hfill
    \begin{subfigure}[b]{0.285\textwidth}
        \centering
        \includegraphics[width=\linewidth]{./figures_experiments/reports/attn_mlp_analysis/attn_mlp_8b_two_hop/expression_region/attn_mlp_expression_region_scores.png}
        \caption{\textit{Qwen3-8B}: Expression}
    \end{subfigure}
    \hfill
    \begin{subfigure}[b]{0.285\textwidth}
        \centering
        \includegraphics[width=\linewidth]{./figures_experiments/reports/attn_mlp_analysis/attn_mlp_8b_two_hop/query_region/attn_mlp_query_region_scores.png}
        \caption{\textit{Qwen3-8B}: Query}
    \end{subfigure}

    \vspace{0.2em}

    \begin{subfigure}[b]{0.285\textwidth}
        \centering
        \includegraphics[width=\linewidth]{./figures_experiments/reports/attn_mlp_analysis/attn_mlp_14b_two_hop/facts_region/attn_mlp_facts_region_scores.png}
        \caption{\textit{Qwen3-14B}: Facts}
    \end{subfigure}
    \hfill
    \begin{subfigure}[b]{0.285\textwidth}
        \centering
        \includegraphics[width=\linewidth]{./figures_experiments/reports/attn_mlp_analysis/attn_mlp_14b_two_hop/expression_region/attn_mlp_expression_region_scores.png}
        \caption{\textit{Qwen3-14B}: Expression}
    \end{subfigure}
    \hfill
    \begin{subfigure}[b]{0.285\textwidth}
        \centering
        \includegraphics[width=\linewidth]{./figures_experiments/reports/attn_mlp_analysis/attn_mlp_14b_two_hop/query_region/attn_mlp_query_region_scores.png}
        \caption{\textit{Qwen3-14B}: Query}
    \end{subfigure}
    \caption{\textbf{Qwen3 facts-first joint attention+MLP galleries for two-hop reasoning.} The joint branch remains closer to the original staged decomposition than pure attention. Facts keep the strongest early-layer dominance, Expression spans the early-to-middle transition, and Query widens mainly in middle and late layers. The same ordering appears in both scales, which supports the view that routing and nonlinear transformation cooperate rather than compete in the joint intervention.}
    \label{fig:qwen_attn_mlp_twohop_facts_gallery}
\end{figure*}

%\begin{figure*}[p]
%    \centering
%    \begin{subfigure}[b]{0.29\textwidth}
%        \centering
%        \includegraphics[width=\linewidth]{./figures_experiments/reports/mlp_analysis/comparison_14B/band_metrics_panel.png}
%        \caption{MLP, 14B \texttt{facts\_first}}
%    \end{subfigure}
%    \hfill
%    \begin{subfigure}[b]{0.29\textwidth}
%        \centering
%        \includegraphics[width=\linewidth]{./figures_experiments/reports/attn_analysis/comparison_14B/band_metrics_panel.png}
%        \caption{Attention, 14B \texttt{facts\_first}}
%    \end{subfigure}
%    \hfill
%    \begin{subfigure}[b]{0.29\textwidth}
%        \centering
%        \includegraphics[width=\linewidth]{./figures_experiments/reports/attn_mlp_analysis/comparison_14B/band_metrics_panel.png}
%        \caption{Attention+MLP, 14B \texttt{facts\_first}}
%    \end{subfigure}
%
%    \vspace{0.2em}
%
%    \begin{subfigure}[b]{0.44\textwidth}
%        \centering
%        \includegraphics[width=\linewidth]{./figures_experiments/reports/attn_analysis_expr_first/comparison_expr_first_14B/band_metrics_panel.png}
%        \caption{Attention, 14B \texttt{expr\_first}}
%    \end{subfigure}
%    \hfill
%    \begin{subfigure}[b]{0.44\textwidth}
%        \centering
%        \includegraphics[width=\linewidth]{./figures_experiments/reports/attn_mlp_analysis_expr_first/comparison_expr_first_14B/band_metrics_panel.png}
%        \caption{Attention+MLP, 14B \texttt{expr\_first}}
%    \end{subfigure}
%    \caption{\textbf{Remaining Qwen3-14B branch-level band summaries.} These panels complete the set of 14B branch summaries across prompt orders. They show the same main tendency already established elsewhere: \texttt{expr\_first} moves more mass into early Expression processing, while the attention branch remains the most middle-heavy and the joint branch the most stage-like.}
%    \label{fig:qwen14_remaining_band_panels}
%\end{figure*}

%\begin{figure*}[p]
%    \centering
%    \begin{subfigure}[b]{0.44\textwidth}
%        \centering
%        \includegraphics[width=\linewidth]{./figures_experiments/reports/attn_analysis_expr_first/comparison_expr_first/band_metrics_panel.png}
%        \caption{Attention, 8B \texttt{expr\_first}}
%    \end{subfigure}
%    \hfill
%    \begin{subfigure}[b]{0.44\textwidth}
%        \centering
%        \includegraphics[width=\linewidth]{./figures_experiments/reports/attn_mlp_analysis_expr_first/comparison_expr_first/band_metrics_panel.png}
%        \caption{Attention+MLP, 8B \texttt{expr\_first}}
%    \end{subfigure}
%    \caption{\textbf{Remaining Qwen3-8B \texttt{expr\_first} branch summaries.} These two omitted panels complete the 8B prompt-order intervention picture for the non-MLP branches. They again show that attention makes the prompt-order shift more structural, while the joint branch retains a cleaner staged decomposition.}
%    \label{fig:qwen8_expr_remaining_band_panels}
%\end{figure*}

%\begin{figure*}[t]
%    \centering
%    \begin{subfigure}[b]{0.44\textwidth}
%        \centering
%        \includegraphics[width=\linewidth]{./figures_experiments/reports/heads_analysis_qwen3_8b_facts_first/validation/plots/pd_abs_ratio_curve.png}
%        \caption{\textit{Qwen3-8B}, \texttt{facts\_first}}
%    \end{subfigure}
%    \hfill
%    \begin{subfigure}[b]{0.44\textwidth}
%        \centering
%        \includegraphics[width=\linewidth]{./figures_experiments/reports/heads_analysis_qwen3_14b_facts_first/validation/plots/pd_abs_ratio_curve.png}
%        \caption{\textit{Qwen3-14B}, \texttt{facts\_first}}
%    \end{subfigure}
%    \caption{\textbf{Qwen3 facts-first absolute validation curves.} These are the two remaining absolute-ratio head-validation plots not previously shown. They complete the Qwen3 Step~3 atlas by pairing the conventional absolute-ratio view with the dPD-shift diagnostics under \texttt{facts\_first}. Both scales preserve the same ordering: Fact-Retrieval is the most damaging single-role ablation, Transmission is second, Splitting is weakest, and \textit{Random-Outside} remains far less harmful.}
%    \label{fig:qwen_heads_abs_ratio_facts}
%\end{figure*}

\begin{figure*}[t]
    \centering
    \begin{subfigure}[b]{0.285\textwidth}
        \centering
        \includegraphics[width=\linewidth]{./figures_experiments/reports/mlp_analysis/mlp_8b_two_hop/facts_region/mlp_facts_region_scores.png}
        \caption{Facts}
    \end{subfigure}
    \hfill
    \begin{subfigure}[b]{0.285\textwidth}
        \centering
        \includegraphics[width=\linewidth]{./figures_experiments/reports/mlp_analysis/mlp_8b_two_hop/expression_region/mlp_expression_region_scores.png}
        \caption{Expression}
    \end{subfigure}
    \hfill
    \begin{subfigure}[b]{0.285\textwidth}
        \centering
        \includegraphics[width=\linewidth]{./figures_experiments/reports/mlp_analysis/mlp_8b_two_hop/query_region/mlp_query_region_scores.png}
        \caption{Query}
    \end{subfigure}
    \caption{\textbf{Qwen3-8B facts-first MLP gallery for two-hop reasoning.} These three curves complete the 8B facts-first MLP atlas. Compared with one-hop, two-hop reasoning weakens the sharpness of every region but preserves the same order: facts dominate earliest, expression remains concentrated in the middle of the network, and query stays strongest in the middle-to-late transition.}
    \label{fig:qwen8_mlp_twohop_facts_gallery}
\end{figure*}

%\begin{figure*}[p]
%    \centering
%    \begin{subfigure}[b]{0.285\textwidth}
%        \centering
%        \includegraphics[width=\linewidth]{./figures_experiments/reports/mlp_analysis/mlp_8b_one_hop/facts_region/mlp_facts_region_abs_dpd_scores.png}
%        \caption{Facts (One-hop)}
%    \end{subfigure}
%    \hfill
%    \begin{subfigure}[b]{0.285\textwidth}
%        \centering
%        \includegraphics[width=\linewidth]{./figures_experiments/reports/mlp_analysis/mlp_8b_one_hop/expression_region/mlp_expression_region_abs_dpd_scores.png}
%        \caption{Expression (One-hop)}
%    \end{subfigure}
%    \hfill
%    \begin{subfigure}[b]{0.285\textwidth}
%        \centering
%        \includegraphics[width=\linewidth]{./figures_experiments/reports/mlp_analysis/mlp_8b_one_hop/query_region/mlp_query_region_abs_dpd_scores.png}
%        \caption{Query (One-hop)}
%    \end{subfigure}
%
%    \vspace{0.2em}
%
%    \begin{subfigure}[b]{0.285\textwidth}
%        \centering
%        \includegraphics[width=\linewidth]{./figures_experiments/reports/mlp_analysis/mlp_8b_two_hop/facts_region/mlp_facts_region_abs_dpd_scores.png}
%        \caption{Facts (Two-hop)}
%    \end{subfigure}
%    \hfill
%    \begin{subfigure}[b]{0.285\textwidth}
%        \centering
%        \includegraphics[width=\linewidth]{./figures_experiments/reports/mlp_analysis/mlp_8b_two_hop/expression_region/mlp_expression_region_abs_dpd_scores.png}
%        \caption{Expression (Two-hop)}
%    \end{subfigure}
%    \hfill
%    \begin{subfigure}[b]{0.285\textwidth}
%        \centering
%        \includegraphics[width=\linewidth]{./figures_experiments/reports/mlp_analysis/mlp_8b_two_hop/query_region/mlp_query_region_abs_dpd_scores.png}
%        \caption{Query (Two-hop)}
%    \end{subfigure}
%    \caption{\textbf{Qwen3-8B facts-first absolute-dPD MLP gallery.} Using $|\mathrm{dPD}|$ instead of signed dPD shows that the same regions dominate in pure magnitude. The two-hop curves are less sharp than one-hop, but the largest-magnitude intervention sites still follow the same facts-to-expression-to-query ordering.}
%    \label{fig:qwen8_mlp_abs_facts_gallery}
%\end{figure*}

%\begin{figure*}[p]
%    \centering
%    \begin{subfigure}[b]{0.285\textwidth}
%        \centering
%        \includegraphics[width=\linewidth]{./figures_experiments/reports/mlp_analysis/mlp_14b_one_hop/facts_region/mlp_facts_region_abs_dpd_scores.png}
%        \caption{Facts (One-hop)}
%    \end{subfigure}
%    \hfill
%    \begin{subfigure}[b]{0.285\textwidth}
%        \centering
%        \includegraphics[width=\linewidth]{./figures_experiments/reports/mlp_analysis/mlp_14b_one_hop/expression_region/mlp_expression_region_abs_dpd_scores.png}
%        \caption{Expression (One-hop)}
%    \end{subfigure}
%    \hfill
%    \begin{subfigure}[b]{0.285\textwidth}
%        \centering
%        \includegraphics[width=\linewidth]{./figures_experiments/reports/mlp_analysis/mlp_14b_one_hop/query_region/mlp_query_region_abs_dpd_scores.png}
%        \caption{Query (One-hop)}
%    \end{subfigure}
%
%    \vspace{0.2em}
%
%    \begin{subfigure}[b]{0.285\textwidth}
%        \centering
%        \includegraphics[width=\linewidth]{./figures_experiments/reports/mlp_analysis/mlp_14b_two_hop/facts_region/mlp_facts_region_abs_dpd_scores.png}
%        \caption{Facts (Two-hop)}
%    \end{subfigure}
%    \hfill
%    \begin{subfigure}[b]{0.285\textwidth}
%        \centering
%        \includegraphics[width=\linewidth]{./figures_experiments/reports/mlp_analysis/mlp_14b_two_hop/expression_region/mlp_expression_region_abs_dpd_scores.png}
%        \caption{Expression (Two-hop)}
%    \end{subfigure}
%    \hfill
%    \begin{subfigure}[b]{0.285\textwidth}
%        \centering
%        \includegraphics[width=\linewidth]{./figures_experiments/reports/mlp_analysis/mlp_14b_two_hop/query_region/mlp_query_region_abs_dpd_scores.png}
%        \caption{Query (Two-hop)}
%    \end{subfigure}
%    \caption{\textbf{Qwen3-14B facts-first absolute-dPD MLP gallery.} The larger model preserves the same magnitude profile as 8B while smoothing some of the late-layer noise. Facts remain the largest early-layer target, Expression retains the clearest middle-layer peak, and Query broadens into the middle-to-late integration stage.}
%    \label{fig:qwen14_mlp_abs_facts_gallery}
%\end{figure*}

%\begin{figure*}[p]
%    \centering
%    \begin{subfigure}[b]{0.285\textwidth}
%        \centering
%        \includegraphics[width=\linewidth]{./figures_experiments/reports/mlp_analysis_expr_first/mlp_8b_one_hop_expr_first/facts_region/mlp_facts_region_abs_dpd_scores.png}
%        \caption{Facts (One-hop)}
%    \end{subfigure}
%    \hfill
%    \begin{subfigure}[b]{0.285\textwidth}
%        \centering
%        \includegraphics[width=\linewidth]{./figures_experiments/reports/mlp_analysis_expr_first/mlp_8b_one_hop_expr_first/expression_region/mlp_expression_region_abs_dpd_scores.png}
%        \caption{Expression (One-hop)}
%    \end{subfigure}
%    \hfill
%    \begin{subfigure}[b]{0.285\textwidth}
%        \centering
%        \includegraphics[width=\linewidth]{./figures_experiments/reports/mlp_analysis_expr_first/mlp_8b_one_hop_expr_first/query_region/mlp_query_region_abs_dpd_scores.png}
%        \caption{Query (One-hop)}
%    \end{subfigure}
%
%    \vspace{0.2em}
%
%    \begin{subfigure}[b]{0.285\textwidth}
%        \centering
%        \includegraphics[width=\linewidth]{./figures_experiments/reports/mlp_analysis_expr_first/mlp_8b_two_hop_expr_first/facts_region/mlp_facts_region_abs_dpd_scores.png}
%        \caption{Facts (Two-hop)}
%    \end{subfigure}
%    \hfill
%    \begin{subfigure}[b]{0.285\textwidth}
%        \centering
%        \includegraphics[width=\linewidth]{./figures_experiments/reports/mlp_analysis_expr_first/mlp_8b_two_hop_expr_first/expression_region/mlp_expression_region_abs_dpd_scores.png}
%        \caption{Expression (Two-hop)}
%    \end{subfigure}
%    \hfill
%    \begin{subfigure}[b]{0.285\textwidth}
%        \centering
%        \includegraphics[width=\linewidth]{./figures_experiments/reports/mlp_analysis_expr_first/mlp_8b_two_hop_expr_first/query_region/mlp_query_region_abs_dpd_scores.png}
%        \caption{Query (Two-hop)}
%    \end{subfigure}
%    \caption{\textbf{Qwen3-8B expr-first absolute-dPD MLP gallery.} In magnitude space, prompt reordering still pulls Expression earlier and broadens Query. The absolute plots show that the prompt-order shift is not only a signed-direction effect: the same front-loaded expression bias is visible even after removing sign.}
%    \label{fig:qwen8_mlp_abs_expr_gallery}
%\end{figure*}

%\begin{figure*}[p]
%    \centering
%    \begin{subfigure}[b]{0.285\textwidth}
%        \centering
%        \includegraphics[width=\linewidth]{./figures_experiments/reports/mlp_analysis_expr_first/mlp_14b_one_hop_expr_first/facts_region/mlp_facts_region_abs_dpd_scores.png}
%        \caption{Facts (One-hop)}
%    \end{subfigure}
%    \hfill
%    \begin{subfigure}[b]{0.285\textwidth}
%        \centering
%        \includegraphics[width=\linewidth]{./figures_experiments/reports/mlp_analysis_expr_first/mlp_14b_one_hop_expr_first/expression_region/mlp_expression_region_abs_dpd_scores.png}
%        \caption{Expression (One-hop)}
%    \end{subfigure}
%    \hfill
%    \begin{subfigure}[b]{0.285\textwidth}
%        \centering
%        \includegraphics[width=\linewidth]{./figures_experiments/reports/mlp_analysis_expr_first/mlp_14b_one_hop_expr_first/query_region/mlp_query_region_abs_dpd_scores.png}
%        \caption{Query (One-hop)}
%    \end{subfigure}
%
%    \vspace{0.2em}
%
%    \begin{subfigure}[b]{0.285\textwidth}
%        \centering
%        \includegraphics[width=\linewidth]{./figures_experiments/reports/mlp_analysis_expr_first/mlp_14b_two_hop_expr_first/facts_region/mlp_facts_region_abs_dpd_scores.png}
%        \caption{Facts (Two-hop)}
%    \end{subfigure}
%    \hfill
%    \begin{subfigure}[b]{0.285\textwidth}
%        \centering
%        \includegraphics[width=\linewidth]{./figures_experiments/reports/mlp_analysis_expr_first/mlp_14b_two_hop_expr_first/expression_region/mlp_expression_region_abs_dpd_scores.png}
%        \caption{Expression (Two-hop)}
%    \end{subfigure}
%    \hfill
%    \begin{subfigure}[b]{0.285\textwidth}
%        \centering
%        \includegraphics[width=\linewidth]{./figures_experiments/reports/mlp_analysis_expr_first/mlp_14b_two_hop_expr_first/query_region/mlp_query_region_abs_dpd_scores.png}
%        \caption{Query (Two-hop)}
%    \end{subfigure}
%    \caption{\textbf{Qwen3-14B expr-first absolute-dPD MLP gallery.} The 14B absolute plots are the cleanest version of the prompt-order intervention at region level. Expression is the dominant early magnitude peak, Facts broaden into the early-to-middle transition, and Query remains a distributed middle-to-late sink.}
%    \label{fig:qwen14_mlp_abs_expr_gallery}
%\end{figure*}

%\begin{figure*}[p]
%    \centering
%    \begin{subfigure}[b]{0.44\textwidth}
%        \centering
%        \includegraphics[width=\linewidth]{./figures_experiments/reports/token_analysis/8b_facts_first_all_hop_raw/category_comparison_simple.png}
%        \caption{\textit{Qwen3-8B}, \texttt{facts\_first}}
%    \end{subfigure}
%    \hfill
%    \begin{subfigure}[b]{0.44\textwidth}
%        \centering
%        \includegraphics[width=\linewidth]{./figures_experiments/reports/token_analysis/14b_facts_first_all_hop_raw/category_comparison_simple.png}
%        \caption{\textit{Qwen3-14B}, \texttt{facts\_first}}
%    \end{subfigure}
%
%    \vspace{0.2em}
%
%    \begin{subfigure}[b]{0.44\textwidth}
%        \centering
%        \includegraphics[width=\linewidth]{./figures_experiments/reports/token_analysis/8b_expr_first_all_hop_raw/category_comparison_simple.png}
%        \caption{\textit{Qwen3-8B}, \texttt{expr\_first}}
%    \end{subfigure}
%    \hfill
%    \begin{subfigure}[b]{0.44\textwidth}
%        \centering
%        \includegraphics[width=\linewidth]{./figures_experiments/reports/token_analysis/14b_expr_first_all_hop_raw/category_comparison_simple.png}
%        \caption{\textit{Qwen3-14B}, \texttt{expr\_first}}
%    \end{subfigure}
%    \caption{\textbf{Qwen3 simple token category comparisons.} These pre-refinement summaries already expose the same core ordering as the refined plots. \texttt{facts\_first} is more fact-heavy, whereas \texttt{expr\_first} shifts visible mass toward query-facing and derived-assignment behavior.}
%    \label{fig:qwen_token_simple_category}
%\end{figure*}

%\begin{figure*}[p]
%    \centering
%    \begin{subfigure}[b]{0.44\textwidth}
%        \centering
%        \includegraphics[width=\linewidth]{./figures_experiments/reports/token_analysis/8b_facts_first_all_hop_raw/layer_stage_simple.png}
%        \caption{\textit{Qwen3-8B}, \texttt{facts\_first}}
%    \end{subfigure}
%    \hfill
%    \begin{subfigure}[b]{0.44\textwidth}
%        \centering
%        \includegraphics[width=\linewidth]{./figures_experiments/reports/token_analysis/14b_facts_first_all_hop_raw/layer_stage_simple.png}
%        \caption{\textit{Qwen3-14B}, \texttt{facts\_first}}
%    \end{subfigure}
%
%    \vspace{0.2em}
%
%    \begin{subfigure}[b]{0.44\textwidth}
%        \centering
%        \includegraphics[width=\linewidth]{./figures_experiments/reports/token_analysis/8b_expr_first_all_hop_raw/layer_stage_simple.png}
%        \caption{\textit{Qwen3-8B}, \texttt{expr\_first}}
%    \end{subfigure}
%    \hfill
%    \begin{subfigure}[b]{0.44\textwidth}
%        \centering
%        \includegraphics[width=\linewidth]{./figures_experiments/reports/token_analysis/14b_expr_first_all_hop_raw/layer_stage_simple.png}
%        \caption{\textit{Qwen3-14B}, \texttt{expr\_first}}
%    \end{subfigure}
%    \caption{\textbf{Qwen3 simple token stage summaries.} The stage-pooled simple plots are a lighter-weight version of the raw heatmaps. They show that the earliest stages remain fact-heavy, while the later stages move toward query-facing categories, with prompt reordering mostly changing the balance rather than the coarse stage layout.}
%    \label{fig:qwen_token_simple_stage}
%\end{figure*}

\paragraph{Cross-model region atlases.}
The cross-model region figures are primarily full attention and joint-region galleries for Llama and Mistral. The summary panels already show that both families place more weight on the Query Region than Qwen3; the layer-wise curves show how that difference is realized in depth. In Llama, query-facing mass is broad and diffuse. In Mistral, the stage ordering is easier to recognize, but query-side routing still occupies much more of the back half of the network than it does in Qwen3. Figures~\ref{fig:llama_attn_facts_gallery}, \ref{fig:llama_attn_expr_gallery}, \ref{fig:llama_attn_mlp_facts_gallery}, and \ref{fig:llama_attn_mlp_expr_gallery} collect the retained Llama attention and joint views. Figures~\ref{fig:mistral_attn_facts_gallery}, \ref{fig:mistral_attn_expr_gallery}, \ref{fig:mistral_attn_mlp_facts_gallery}, and \ref{fig:mistral_attn_mlp_expr_gallery} provide the corresponding Mistral views.

\begin{figure*}[t]
    \centering
    \begin{subfigure}[b]{0.285\textwidth}
        \centering
        \includegraphics[width=\linewidth]{./figures_experiments/reports_llama3/attn_analysis_facts_first/attn_one_hop_facts_first/facts_region/attn_facts_region_scores.png}
        \caption{Facts (One-hop)}
    \end{subfigure}
    \hfill
    \begin{subfigure}[b]{0.285\textwidth}
        \centering
        \includegraphics[width=\linewidth]{./figures_experiments/reports_llama3/attn_analysis_facts_first/attn_one_hop_facts_first/expression_region/attn_expression_region_scores.png}
        \caption{Expression (One-hop)}
    \end{subfigure}
    \hfill
    \begin{subfigure}[b]{0.285\textwidth}
        \centering
        \includegraphics[width=\linewidth]{./figures_experiments/reports_llama3/attn_analysis_facts_first/attn_one_hop_facts_first/query_region/attn_query_region_scores.png}
        \caption{Query (One-hop)}
    \end{subfigure}

    \vspace{0.2em}

    \begin{subfigure}[b]{0.285\textwidth}
        \centering
        \includegraphics[width=\linewidth]{./figures_experiments/reports_llama3/attn_analysis_facts_first/attn_two_hop_facts_first/facts_region/attn_facts_region_scores.png}
        \caption{Facts (Two-hop)}
    \end{subfigure}
    \hfill
    \begin{subfigure}[b]{0.285\textwidth}
        \centering
        \includegraphics[width=\linewidth]{./figures_experiments/reports_llama3/attn_analysis_facts_first/attn_two_hop_facts_first/expression_region/attn_expression_region_scores.png}
        \caption{Expression (Two-hop)}
    \end{subfigure}
    \hfill
    \begin{subfigure}[b]{0.285\textwidth}
        \centering
        \includegraphics[width=\linewidth]{./figures_experiments/reports_llama3/attn_analysis_facts_first/attn_two_hop_facts_first/query_region/attn_query_region_scores.png}
        \caption{Query (Two-hop)}
    \end{subfigure}
    \caption{\textbf{Llama attention-region gallery under \texttt{facts\_first}.} Llama's attention branch is strongly query-heavy even without prompt reordering. Facts and expression remain early dominant, but the Query Region keeps unusually large middle and late mass in both one-hop and two-hop reasoning.}
    \label{fig:llama_attn_facts_gallery}
\end{figure*}

\begin{figure*}[t]
    \centering
    \begin{subfigure}[b]{0.31\textwidth}
        \centering
        \includegraphics[width=\linewidth]{./figures_experiments/reports_llama3/attn_analysis_expr_first/attn_one_hop_expr_first/facts_region/attn_facts_region_scores.png}
        \caption{Facts (One-hop)}
    \end{subfigure}
    \hfill
    \begin{subfigure}[b]{0.31\textwidth}
        \centering
        \includegraphics[width=\linewidth]{./figures_experiments/reports_llama3/attn_analysis_expr_first/attn_one_hop_expr_first/expression_region/attn_expression_region_scores.png}
        \caption{Expression (One-hop)}
    \end{subfigure}
    \hfill
    \begin{subfigure}[b]{0.31\textwidth}
        \centering
        \includegraphics[width=\linewidth]{./figures_experiments/reports_llama3/attn_analysis_expr_first/attn_one_hop_expr_first/query_region/attn_query_region_scores.png}
        \caption{Query (One-hop)}
    \end{subfigure}

    \vspace{0.15em}

    \begin{subfigure}[b]{0.31\textwidth}
        \centering
        \includegraphics[width=\linewidth]{./figures_experiments/reports_llama3/attn_analysis_expr_first/attn_two_hop_expr_first/facts_region/attn_facts_region_scores.png}
        \caption{Facts (Two-hop)}
    \end{subfigure}
    \hfill
    \begin{subfigure}[b]{0.31\textwidth}
        \centering
        \includegraphics[width=\linewidth]{./figures_experiments/reports_llama3/attn_analysis_expr_first/attn_two_hop_expr_first/expression_region/attn_expression_region_scores.png}
        \caption{Expression (Two-hop)}
    \end{subfigure}
    \hfill
    \begin{subfigure}[b]{0.31\textwidth}
        \centering
        \includegraphics[width=\linewidth]{./figures_experiments/reports_llama3/attn_analysis_expr_first/attn_two_hop_expr_first/query_region/attn_query_region_scores.png}
        \caption{Query (Two-hop)}
    \end{subfigure}
    \caption{\textbf{Llama attention-region gallery under \texttt{expr\_first}.} Front-loading the expression pushes more mass into early Expression processing, but the Query Region remains broad in both one-hop and two-hop reasoning. The two-hop row further widens the query-facing stage, which makes Llama's diffuse, query-centered routing especially clear.}
    \label{fig:llama_attn_expr_gallery}
\end{figure*}

\begin{figure*}[t]
    \centering
    \begin{subfigure}[b]{0.31\textwidth}
        \centering
        \includegraphics[width=\linewidth]{./figures_experiments/reports_llama3/attn_mlp_analysis_facts_first/attn_mlp_one_hop_facts_first/facts_region/attn_mlp_facts_region_scores.png}
        \caption{Facts (One-hop)}
    \end{subfigure}
    \hfill
    \begin{subfigure}[b]{0.31\textwidth}
        \centering
        \includegraphics[width=\linewidth]{./figures_experiments/reports_llama3/attn_mlp_analysis_facts_first/attn_mlp_one_hop_facts_first/expression_region/attn_mlp_expression_region_scores.png}
        \caption{Expression (One-hop)}
    \end{subfigure}
    \hfill
    \begin{subfigure}[b]{0.31\textwidth}
        \centering
        \includegraphics[width=\linewidth]{./figures_experiments/reports_llama3/attn_mlp_analysis_facts_first/attn_mlp_one_hop_facts_first/query_region/attn_mlp_query_region_scores.png}
        \caption{Query (One-hop)}
    \end{subfigure}

    \vspace{0.15em}

    \begin{subfigure}[b]{0.31\textwidth}
        \centering
        \includegraphics[width=\linewidth]{./figures_experiments/reports_llama3/attn_mlp_analysis_facts_first/attn_mlp_two_hop_facts_first/facts_region/attn_mlp_facts_region_scores.png}
        \caption{Facts (Two-hop)}
    \end{subfigure}
    \hfill
    \begin{subfigure}[b]{0.31\textwidth}
        \centering
        \includegraphics[width=\linewidth]{./figures_experiments/reports_llama3/attn_mlp_analysis_facts_first/attn_mlp_two_hop_facts_first/expression_region/attn_mlp_expression_region_scores.png}
        \caption{Expression (Two-hop)}
    \end{subfigure}
    \hfill
    \begin{subfigure}[b]{0.31\textwidth}
        \centering
        \includegraphics[width=\linewidth]{./figures_experiments/reports_llama3/attn_mlp_analysis_facts_first/attn_mlp_two_hop_facts_first/query_region/attn_mlp_query_region_scores.png}
        \caption{Query (Two-hop)}
    \end{subfigure}
    \caption{\textbf{Llama joint attention+MLP gallery under \texttt{facts\_first}.} The joint branch is somewhat more stage-like than pure attention, but still much more query-heavy than Qwen3. Even in one-hop reasoning the Query Region remains broad across the back half of the network, and the two-hop row broadens this profile further.}
    \label{fig:llama_attn_mlp_facts_gallery}
\end{figure*}

\begin{figure*}[t]
    \centering
    \begin{subfigure}[b]{0.31\textwidth}
        \centering
        \includegraphics[width=\linewidth]{./figures_experiments/reports_llama3/attn_mlp_analysis_expr_first/attn_mlp_one_hop_expr_first/facts_region/attn_mlp_facts_region_scores.png}
        \caption{Facts (One-hop)}
    \end{subfigure}
    \hfill
    \begin{subfigure}[b]{0.31\textwidth}
        \centering
        \includegraphics[width=\linewidth]{./figures_experiments/reports_llama3/attn_mlp_analysis_expr_first/attn_mlp_one_hop_expr_first/expression_region/attn_mlp_expression_region_scores.png}
        \caption{Expression (One-hop)}
    \end{subfigure}
    \hfill
    \begin{subfigure}[b]{0.31\textwidth}
        \centering
        \includegraphics[width=\linewidth]{./figures_experiments/reports_llama3/attn_mlp_analysis_expr_first/attn_mlp_one_hop_expr_first/query_region/attn_mlp_query_region_scores.png}
        \caption{Query (One-hop)}
    \end{subfigure}

    \vspace{0.15em}

    \begin{subfigure}[b]{0.31\textwidth}
        \centering
        \includegraphics[width=\linewidth]{./figures_experiments/reports_llama3/attn_mlp_analysis_expr_first/attn_mlp_two_hop_expr_first/facts_region/attn_mlp_facts_region_scores.png}
        \caption{Facts (Two-hop)}
    \end{subfigure}
    \hfill
    \begin{subfigure}[b]{0.31\textwidth}
        \centering
        \includegraphics[width=\linewidth]{./figures_experiments/reports_llama3/attn_mlp_analysis_expr_first/attn_mlp_two_hop_expr_first/expression_region/attn_mlp_expression_region_scores.png}
        \caption{Expression (Two-hop)}
    \end{subfigure}
    \hfill
    \begin{subfigure}[b]{0.31\textwidth}
        \centering
        \includegraphics[width=\linewidth]{./figures_experiments/reports_llama3/attn_mlp_analysis_expr_first/attn_mlp_two_hop_expr_first/query_region/attn_mlp_query_region_scores.png}
        \caption{Query (Two-hop)}
    \end{subfigure}
    \caption{\textbf{Llama joint attention+MLP gallery under \texttt{expr\_first}.} The expression-first intervention sharpens the early Expression peak, but the Query Region remains broad in both rows. Two-hop reasoning further stretches the query-facing stage across the back half of the network, which keeps the joint branch visibly more diffuse than its Qwen3 counterpart.}
    \label{fig:llama_attn_mlp_expr_gallery}
\end{figure*}

\begin{figure*}[t]
    \centering
    \begin{subfigure}[b]{0.31\textwidth}
        \centering
        \includegraphics[width=\linewidth]{./figures_experiments/reports_mistral/attn_analysis_facts_first/attn_one_hop_facts_first/facts_region/attn_facts_region_scores.png}
        \caption{Facts (One-hop)}
    \end{subfigure}
    \hfill
    \begin{subfigure}[b]{0.31\textwidth}
        \centering
        \includegraphics[width=\linewidth]{./figures_experiments/reports_mistral/attn_analysis_facts_first/attn_one_hop_facts_first/expression_region/attn_expression_region_scores.png}
        \caption{Expression (One-hop)}
    \end{subfigure}
    \hfill
    \begin{subfigure}[b]{0.31\textwidth}
        \centering
        \includegraphics[width=\linewidth]{./figures_experiments/reports_mistral/attn_analysis_facts_first/attn_one_hop_facts_first/query_region/attn_query_region_scores.png}
        \caption{Query (One-hop)}
    \end{subfigure}

    \vspace{0.15em}

    \begin{subfigure}[b]{0.31\textwidth}
        \centering
        \includegraphics[width=\linewidth]{./figures_experiments/reports_mistral/attn_analysis_facts_first/attn_two_hop_facts_first/facts_region/attn_facts_region_scores.png}
        \caption{Facts (Two-hop)}
    \end{subfigure}
    \hfill
    \begin{subfigure}[b]{0.31\textwidth}
        \centering
        \includegraphics[width=\linewidth]{./figures_experiments/reports_mistral/attn_analysis_facts_first/attn_two_hop_facts_first/expression_region/attn_expression_region_scores.png}
        \caption{Expression (Two-hop)}
    \end{subfigure}
    \hfill
    \begin{subfigure}[b]{0.31\textwidth}
        \centering
        \includegraphics[width=\linewidth]{./figures_experiments/reports_mistral/attn_analysis_facts_first/attn_two_hop_facts_first/query_region/attn_query_region_scores.png}
        \caption{Query (Two-hop)}
    \end{subfigure}
    \caption{\textbf{Mistral attention-region gallery under \texttt{facts\_first}.} Mistral retains recognizable early Facts and Expression processing in both rows, but keeps considerably more query mass than Qwen3. The two-hop row broadens this query-facing profile further across the middle and late layers.}
    \label{fig:mistral_attn_facts_gallery}
\end{figure*}

\begin{figure*}[t]
    \centering
    \begin{subfigure}[b]{0.285\textwidth}
        \centering
        \includegraphics[width=\linewidth]{./figures_experiments/reports_mistral/attn_analysis_expr_first/attn_one_hop_expr_first/facts_region/attn_facts_region_scores.png}
        \caption{Facts (One-hop)}
    \end{subfigure}
    \hfill
    \begin{subfigure}[b]{0.285\textwidth}
        \centering
        \includegraphics[width=\linewidth]{./figures_experiments/reports_mistral/attn_analysis_expr_first/attn_one_hop_expr_first/expression_region/attn_expression_region_scores.png}
        \caption{Expression (One-hop)}
    \end{subfigure}
    \hfill
    \begin{subfigure}[b]{0.285\textwidth}
        \centering
        \includegraphics[width=\linewidth]{./figures_experiments/reports_mistral/attn_analysis_expr_first/attn_one_hop_expr_first/query_region/attn_query_region_scores.png}
        \caption{Query (One-hop)}
    \end{subfigure}

    \vspace{0.2em}

    \begin{subfigure}[b]{0.285\textwidth}
        \centering
        \includegraphics[width=\linewidth]{./figures_experiments/reports_mistral/attn_analysis_expr_first/attn_two_hop_expr_first/facts_region/attn_facts_region_scores.png}
        \caption{Facts (Two-hop)}
    \end{subfigure}
    \hfill
    \begin{subfigure}[b]{0.285\textwidth}
        \centering
        \includegraphics[width=\linewidth]{./figures_experiments/reports_mistral/attn_analysis_expr_first/attn_two_hop_expr_first/expression_region/attn_expression_region_scores.png}
        \caption{Expression (Two-hop)}
    \end{subfigure}
    \hfill
    \begin{subfigure}[b]{0.285\textwidth}
        \centering
        \includegraphics[width=\linewidth]{./figures_experiments/reports_mistral/attn_analysis_expr_first/attn_two_hop_expr_first/query_region/attn_query_region_scores.png}
        \caption{Query (Two-hop)}
    \end{subfigure}
    \caption{\textbf{Mistral attention-region gallery under \texttt{expr\_first}.} Prompt reordering shifts more weight into early Expression processing but leaves the Query Region broad and comparatively heavy. The accompanying band summary below confirms that Mistral's attention branch remains highly query-facing even under \texttt{expr\_first}.}
    \label{fig:mistral_attn_expr_gallery}
\end{figure*}

%\begin{figure*}[p]
%    \centering
%    \begin{subfigure}[b]{0.44\textwidth}
%        \centering
%        \includegraphics[width=\linewidth]{./figures_experiments/reports_mistral/attn_analysis_expr_first/comparison/band_metrics_panel.png}
%        \caption{Attention branch}
%    \end{subfigure}
%    \hfill
%    \begin{subfigure}[b]{0.44\textwidth}
%        \centering
%        \includegraphics[width=\linewidth]{./figures_experiments/reports_mistral/attn_mlp_analysis_expr_first/comparison/band_metrics_panel.png}
%        \caption{Joint attention+MLP branch}
%    \end{subfigure}
%    \caption{\textbf{Remaining Mistral expr-first branch summaries.} These two band panels complete the Mistral prompt-order atlas for the non-MLP branches. The attention panel is more middle- and query-heavy; the joint panel is somewhat more stage-like but still broader than its Qwen3 counterpart.}
%    \label{fig:mistral_expr_remaining_band_panels}
%\end{figure*}

%\begin{figure*}[p]
%    \centering
%    \includegraphics[width=\textwidth]{./figures_experiments/reports_mistral/attn_mlp_analysis_facts_first/comparison/band_metrics_panel.png}
%    \caption{\textbf{Mistral facts-first joint attention+MLP band summary.} This final omitted branch-level summary completes the Mistral band-panel atlas. The joint intervention preserves the same broad staged template as the MLP branch, but still places comparatively strong mass on the Query Region throughout the back half of the network.}
%    \label{fig:mistral_attn_mlp_facts_band_panel}
%\end{figure*}

\begin{figure*}[t]
    \centering
    \begin{subfigure}[b]{0.285\textwidth}
        \centering
        \includegraphics[width=\linewidth]{./figures_experiments/reports_mistral/attn_mlp_analysis_facts_first/attn_mlp_one_hop_facts_first/facts_region/attn_mlp_facts_region_scores.png}
        \caption{Facts (One-hop)}
    \end{subfigure}
    \hfill
    \begin{subfigure}[b]{0.285\textwidth}
        \centering
        \includegraphics[width=\linewidth]{./figures_experiments/reports_mistral/attn_mlp_analysis_facts_first/attn_mlp_one_hop_facts_first/expression_region/attn_mlp_expression_region_scores.png}
        \caption{Expression (One-hop)}
    \end{subfigure}
    \hfill
    \begin{subfigure}[b]{0.285\textwidth}
        \centering
        \includegraphics[width=\linewidth]{./figures_experiments/reports_mistral/attn_mlp_analysis_facts_first/attn_mlp_one_hop_facts_first/query_region/attn_mlp_query_region_scores.png}
        \caption{Query (One-hop)}
    \end{subfigure}

    \vspace{0.2em}

    \begin{subfigure}[b]{0.285\textwidth}
        \centering
        \includegraphics[width=\linewidth]{./figures_experiments/reports_mistral/attn_mlp_analysis_facts_first/attn_mlp_two_hop_facts_first/facts_region/attn_mlp_facts_region_scores.png}
        \caption{Facts (Two-hop)}
    \end{subfigure}
    \hfill
    \begin{subfigure}[b]{0.285\textwidth}
        \centering
        \includegraphics[width=\linewidth]{./figures_experiments/reports_mistral/attn_mlp_analysis_facts_first/attn_mlp_two_hop_facts_first/expression_region/attn_mlp_expression_region_scores.png}
        \caption{Expression (Two-hop)}
    \end{subfigure}
    \hfill
    \begin{subfigure}[b]{0.285\textwidth}
        \centering
        \includegraphics[width=\linewidth]{./figures_experiments/reports_mistral/attn_mlp_analysis_facts_first/attn_mlp_two_hop_facts_first/query_region/attn_mlp_query_region_scores.png}
        \caption{Query (Two-hop)}
    \end{subfigure}
    \caption{\textbf{Mistral joint attention+MLP gallery under \texttt{facts\_first}.} The joint branch again looks more stage-like than pure attention, but still preserves a heavy Query Region relative to Qwen3. The full curves show that Mistral's joint intervention sharpens early Facts while leaving a substantial amount of middle-to-late query-facing mass.}
    \label{fig:mistral_attn_mlp_facts_gallery}
\end{figure*}

\begin{figure*}[t]
    \centering
    \begin{subfigure}[b]{0.31\textwidth}
        \centering
        \includegraphics[width=\linewidth]{./figures_experiments/reports_mistral/attn_mlp_analysis_expr_first/attn_mlp_one_hop_expr_first/facts_region/attn_mlp_facts_region_scores.png}
        \caption{Facts (One-hop)}
    \end{subfigure}
    \hfill
    \begin{subfigure}[b]{0.31\textwidth}
        \centering
        \includegraphics[width=\linewidth]{./figures_experiments/reports_mistral/attn_mlp_analysis_expr_first/attn_mlp_one_hop_expr_first/expression_region/attn_mlp_expression_region_scores.png}
        \caption{Expression (One-hop)}
    \end{subfigure}
    \hfill
    \begin{subfigure}[b]{0.31\textwidth}
        \centering
        \includegraphics[width=\linewidth]{./figures_experiments/reports_mistral/attn_mlp_analysis_expr_first/attn_mlp_one_hop_expr_first/query_region/attn_mlp_query_region_scores.png}
        \caption{Query (One-hop)}
    \end{subfigure}

    \vspace{0.15em}

    \begin{subfigure}[b]{0.31\textwidth}
        \centering
        \includegraphics[width=\linewidth]{./figures_experiments/reports_mistral/attn_mlp_analysis_expr_first/attn_mlp_two_hop_expr_first/facts_region/attn_mlp_facts_region_scores.png}
        \caption{Facts (Two-hop)}
    \end{subfigure}
    \hfill
    \begin{subfigure}[b]{0.31\textwidth}
        \centering
        \includegraphics[width=\linewidth]{./figures_experiments/reports_mistral/attn_mlp_analysis_expr_first/attn_mlp_two_hop_expr_first/expression_region/attn_mlp_expression_region_scores.png}
        \caption{Expression (Two-hop)}
    \end{subfigure}
    \hfill
    \begin{subfigure}[b]{0.31\textwidth}
        \centering
        \includegraphics[width=\linewidth]{./figures_experiments/reports_mistral/attn_mlp_analysis_expr_first/attn_mlp_two_hop_expr_first/query_region/attn_mlp_query_region_scores.png}
        \caption{Query (Two-hop)}
    \end{subfigure}
    \caption{\textbf{Mistral joint attention+MLP gallery under \texttt{expr\_first}.} In one-hop reasoning, front-loading the expression amplifies the early Expression Region while keeping the Query Region broad. Two-hop reasoning preserves the same reordered stage structure but broadens the Query Region further, so the joint branch remains visibly less cleanly factorized and more query-facing than its Qwen3 counterpart.}
    \label{fig:mistral_attn_mlp_expr_gallery}
\end{figure*}

%\begin{figure*}[p]
%    \centering
%    \begin{subfigure}[b]{0.285\textwidth}
%        \centering
%        \includegraphics[width=\linewidth]{./figures_experiments/reports_llama3/mlp_analysis_facts_first/mlp_one_hop_facts_first/facts_region/mlp_facts_region_abs_dpd_scores.png}
%        \caption{Llama, facts-first (one-hop)}
%    \end{subfigure}
%    \hfill
%    \begin{subfigure}[b]{0.285\textwidth}
%        \centering
%        \includegraphics[width=\linewidth]{./figures_experiments/reports_llama3/mlp_analysis_facts_first/mlp_two_hop_facts_first/facts_region/mlp_facts_region_abs_dpd_scores.png}
%        \caption{Llama, facts-first (two-hop)}
%    \end{subfigure}
%    \hfill
%    \begin{subfigure}[b]{0.285\textwidth}
%        \centering
%        \includegraphics[width=\linewidth]{./figures_experiments/reports_llama3/mlp_analysis_expr_first/mlp_one_hop_expr_first/facts_region/mlp_facts_region_abs_dpd_scores.png}
%        \caption{Llama, expr-first (one-hop)}
%    \end{subfigure}
%
%    \vspace{0.2em}
%
%    \begin{subfigure}[b]{0.285\textwidth}
%        \centering
%        \includegraphics[width=\linewidth]{./figures_experiments/reports_llama3/mlp_analysis_expr_first/mlp_two_hop_expr_first/facts_region/mlp_facts_region_abs_dpd_scores.png}
%        \caption{Llama, expr-first (two-hop)}
%    \end{subfigure}
%    \hfill
%    \begin{subfigure}[b]{0.285\textwidth}
%        \centering
%        \includegraphics[width=\linewidth]{./figures_experiments/reports_mistral/mlp_analysis_facts_first/mlp_one_hop_facts_first/facts_region/mlp_facts_region_abs_dpd_scores.png}
%        \caption{Mistral, facts-first (one-hop)}
%    \end{subfigure}
%    \hfill
%    \begin{subfigure}[b]{0.285\textwidth}
%        \centering
%        \includegraphics[width=\linewidth]{./figures_experiments/reports_mistral/mlp_analysis_facts_first/mlp_two_hop_facts_first/facts_region/mlp_facts_region_abs_dpd_scores.png}
%        \caption{Mistral, facts-first (two-hop)}
%    \end{subfigure}
%    \caption{\textbf{Cross-model absolute-dPD facts-region diagnostics.} The absolute facts-region curves make it clear that both non-Qwen families still retain a meaningful facts-processing stage in magnitude space. The main difference from Qwen3 is therefore not the absence of facts processing, but its weaker separation from the query-facing stages.}
%    \label{fig:cross_model_abs_facts_region}
%\end{figure*}

%\begin{figure*}[p]
%    \centering
%    \begin{subfigure}[b]{0.285\textwidth}
%        \centering
%        \includegraphics[width=\linewidth]{./figures_experiments/reports_mistral/mlp_analysis_expr_first/mlp_one_hop_expr_first/facts_region/mlp_facts_region_abs_dpd_scores.png}
%        \caption{Mistral, expr-first (one-hop)}
%    \end{subfigure}
%    \hfill
%    \begin{subfigure}[b]{0.285\textwidth}
%        \centering
%        \includegraphics[width=\linewidth]{./figures_experiments/reports_mistral/mlp_analysis_expr_first/mlp_two_hop_expr_first/facts_region/mlp_facts_region_abs_dpd_scores.png}
%        \caption{Mistral, expr-first (two-hop)}
%    \end{subfigure}
%    \hfill
%    \begin{subfigure}[b]{0.285\textwidth}
%        \centering
%        \includegraphics[width=\linewidth]{./figures_experiments/reports_llama3/mlp_analysis_facts_first/mlp_one_hop_facts_first/expression_region/mlp_expression_region_abs_dpd_scores.png}
%        \caption{Llama facts-first: Expression}
%    \end{subfigure}
%
%    \vspace{0.2em}
%
%    \begin{subfigure}[b]{0.285\textwidth}
%        \centering
%        \includegraphics[width=\linewidth]{./figures_experiments/reports_llama3/mlp_analysis_expr_first/mlp_one_hop_expr_first/expression_region/mlp_expression_region_abs_dpd_scores.png}
%        \caption{Llama expr-first: Expression}
%    \end{subfigure}
%    \hfill
%    \begin{subfigure}[b]{0.285\textwidth}
%        \centering
%        \includegraphics[width=\linewidth]{./figures_experiments/reports_mistral/mlp_analysis_facts_first/mlp_one_hop_facts_first/expression_region/mlp_expression_region_abs_dpd_scores.png}
%        \caption{Mistral facts-first: Expression}
%    \end{subfigure}
%    \hfill
%    \begin{subfigure}[b]{0.285\textwidth}
%        \centering
%        \includegraphics[width=\linewidth]{./figures_experiments/reports_mistral/mlp_analysis_expr_first/mlp_one_hop_expr_first/expression_region/mlp_expression_region_abs_dpd_scores.png}
%        \caption{Mistral expr-first: Expression}
%    \end{subfigure}
%    \caption{\textbf{Additional cross-model absolute-dPD diagnostics for facts and expression processing.} These remaining magnitude plots reinforce the same conclusion as the signed and BCR summaries: prompt reordering increases early expression mass in both Llama and Mistral, but neither model develops the same clean fact-versus-query separation seen in Qwen3.}
%    \label{fig:cross_model_abs_remaining}
%\end{figure*}
%
%\begin{figure*}[p]
%    \centering
%    \begin{subfigure}[b]{0.285\textwidth}
%        \centering
%        \includegraphics[width=\linewidth]{./figures_experiments/reports_llama3/mlp_analysis_facts_first/mlp_one_hop_facts_first/query_region/mlp_query_region_abs_dpd_scores.png}
%        \caption{Facts-first: Query (One-hop)}
%    \end{subfigure}
%    \hfill
%    \begin{subfigure}[b]{0.285\textwidth}
%        \centering
%        \includegraphics[width=\linewidth]{./figures_experiments/reports_llama3/mlp_analysis_facts_first/mlp_two_hop_facts_first/expression_region/mlp_expression_region_abs_dpd_scores.png}
%        \caption{Facts-first: Expression (Two-hop)}
%    \end{subfigure}
%    \hfill
%    \begin{subfigure}[b]{0.285\textwidth}
%        \centering
%        \includegraphics[width=\linewidth]{./figures_experiments/reports_llama3/mlp_analysis_facts_first/mlp_two_hop_facts_first/query_region/mlp_query_region_abs_dpd_scores.png}
%        \caption{Facts-first: Query (Two-hop)}
%    \end{subfigure}
%
%    \vspace{0.2em}
%
%    \begin{subfigure}[b]{0.285\textwidth}
%        \centering
%        \includegraphics[width=\linewidth]{./figures_experiments/reports_llama3/mlp_analysis_expr_first/mlp_one_hop_expr_first/query_region/mlp_query_region_abs_dpd_scores.png}
%        \caption{Expr-first: Query (One-hop)}
%    \end{subfigure}
%    \hfill
%    \begin{subfigure}[b]{0.285\textwidth}
%        \centering
%        \includegraphics[width=\linewidth]{./figures_experiments/reports_llama3/mlp_analysis_expr_first/mlp_two_hop_expr_first/expression_region/mlp_expression_region_abs_dpd_scores.png}
%        \caption{Expr-first: Expression (Two-hop)}
%    \end{subfigure}
%    \hfill
%    \begin{subfigure}[b]{0.285\textwidth}
%        \centering
%        \includegraphics[width=\linewidth]{./figures_experiments/reports_llama3/mlp_analysis_expr_first/mlp_two_hop_expr_first/query_region/mlp_query_region_abs_dpd_scores.png}
%        \caption{Expr-first: Query (Two-hop)}
%    \end{subfigure}
%    \caption{\textbf{Residual Llama absolute-dPD diagnostics.} These six remaining magnitude plots complete the Llama MLP atlas. They show that the most persistent residual magnitude sits in the Query Region, especially once the query is moved to the front or when reasoning depth increases to two-hop.}
%    \label{fig:llama_abs_residual_gallery}
%\end{figure*}

%\begin{figure*}[p]
%    \centering
%    \begin{subfigure}[b]{0.285\textwidth}
%        \centering
%        \includegraphics[width=\linewidth]{./figures_experiments/reports_mistral/mlp_analysis_facts_first/mlp_one_hop_facts_first/query_region/mlp_query_region_abs_dpd_scores.png}
%        \caption{Facts-first: Query (One-hop)}
%    \end{subfigure}
%    \hfill
%    \begin{subfigure}[b]{0.285\textwidth}
%        \centering
%        \includegraphics[width=\linewidth]{./figures_experiments/reports_mistral/mlp_analysis_facts_first/mlp_two_hop_facts_first/expression_region/mlp_expression_region_abs_dpd_scores.png}
%        \caption{Facts-first: Expression (Two-hop)}
%    \end{subfigure}
%    \hfill
%    \begin{subfigure}[b]{0.285\textwidth}
%        \centering
%        \includegraphics[width=\linewidth]{./figures_experiments/reports_mistral/mlp_analysis_facts_first/mlp_two_hop_facts_first/query_region/mlp_query_region_abs_dpd_scores.png}
%        \caption{Facts-first: Query (Two-hop)}
%    \end{subfigure}
%
%    \vspace{0.2em}
%
%    \begin{subfigure}[b]{0.285\textwidth}
%        \centering
%        \includegraphics[width=\linewidth]{./figures_experiments/reports_mistral/mlp_analysis_expr_first/mlp_one_hop_expr_first/query_region/mlp_query_region_abs_dpd_scores.png}
%        \caption{Expr-first: Query (One-hop)}
%    \end{subfigure}
%    \hfill
%    \begin{subfigure}[b]{0.285\textwidth}
%        \centering
%        \includegraphics[width=\linewidth]{./figures_experiments/reports_mistral/mlp_analysis_expr_first/mlp_two_hop_expr_first/expression_region/mlp_expression_region_abs_dpd_scores.png}
%        \caption{Expr-first: Expression (Two-hop)}
%    \end{subfigure}
%    \hfill
%    \begin{subfigure}[b]{0.285\textwidth}
%        \centering
%        \includegraphics[width=\linewidth]{./figures_experiments/reports_mistral/mlp_analysis_expr_first/mlp_two_hop_expr_first/query_region/mlp_query_region_abs_dpd_scores.png}
%        \caption{Expr-first: Query (Two-hop)}
%    \end{subfigure}
%    \caption{\textbf{Residual Mistral absolute-dPD diagnostics.} These final magnitude plots complete the Mistral MLP atlas. They reinforce the same picture seen throughout the appendix: Mistral preserves a recognizable expression-processing stage, but its strongest residual magnitude is often query-facing once the task becomes two-hop or the prompt is reordered.}
%    \label{fig:mistral_abs_residual_gallery}
%\end{figure*}

%\begin{figure*}[p]
%    \centering
%    \begin{subfigure}[b]{0.44\textwidth}
%        \centering
%        \includegraphics[width=\linewidth]{./figures_experiments/reports_llama3/token_analysis_facts_first_raw/category_comparison_simple.png}
%        \caption{Llama, \texttt{facts\_first}}
%    \end{subfigure}
%    \hfill
%    \begin{subfigure}[b]{0.44\textwidth}
%        \centering
%        \includegraphics[width=\linewidth]{./figures_experiments/reports_llama3/token_analysis_expr_first_raw/category_comparison_simple.png}
%        \caption{Llama, \texttt{expr\_first}}
%    \end{subfigure}
%
%    \vspace{0.2em}
%
%    \begin{subfigure}[b]{0.44\textwidth}
%        \centering
%        \includegraphics[width=\linewidth]{./figures_experiments/reports_mistral/token_analysis_facts_first_raw/category_comparison_simple.png}
%        \caption{Mistral, \texttt{facts\_first}}
%    \end{subfigure}
%    \hfill
%    \begin{subfigure}[b]{0.44\textwidth}
%        \centering
%        \includegraphics[width=\linewidth]{./figures_experiments/reports_mistral/token_analysis_expr_first_raw/category_comparison_simple.png}
%        \caption{Mistral, \texttt{expr\_first}}
%    \end{subfigure}
%    \caption{\textbf{Cross-model simple token category comparisons.} These non-refined token summaries already show that Llama and Mistral allocate more visible mass to query-facing categories than Qwen3. The refined analyses sharpen the same picture rather than changing it.}
%    \label{fig:cross_model_token_simple_category}
%\end{figure*}

%\begin{figure*}[p]
%    \centering
%    \begin{subfigure}[b]{0.44\textwidth}
%        \centering
%        \includegraphics[width=\linewidth]{./figures_experiments/reports_llama3/token_analysis_facts_first_raw/layer_stage_simple.png}
%        \caption{Llama, \texttt{facts\_first}}
%    \end{subfigure}
%    \hfill
%    \begin{subfigure}[b]{0.44\textwidth}
%        \centering
%        \includegraphics[width=\linewidth]{./figures_experiments/reports_llama3/token_analysis_expr_first_raw/layer_stage_simple.png}
%        \caption{Llama, \texttt{expr\_first}}
%    \end{subfigure}
%
%    \vspace{0.2em}
%
%    \begin{subfigure}[b]{0.44\textwidth}
%        \centering
%        \includegraphics[width=\linewidth]{./figures_experiments/reports_mistral/token_analysis_facts_first_raw/layer_stage_simple.png}
%        \caption{Mistral, \texttt{facts\_first}}
%    \end{subfigure}
%    \hfill
%    \begin{subfigure}[b]{0.44\textwidth}
%        \centering
%        \includegraphics[width=\linewidth]{./figures_experiments/reports_mistral/token_analysis_expr_first_raw/layer_stage_simple.png}
%        \caption{Mistral, \texttt{expr\_first}}
%    \end{subfigure}
%    \caption{\textbf{Cross-model simple token stage summaries.} Stage-pooled raw token plots again show that both non-Qwen models are more query-facing. Llama is the most diffuse, while Mistral preserves a somewhat cleaner early-to-late progression but still places more late-stage weight on query-region categories than Qwen3.}
%    \label{fig:cross_model_token_simple_stage}
%\end{figure*}
